% !TeX program = xelatex
\documentclass[11pt,a4paper]{article}
% BEGIN INCLUDED FILE: preamble.tex
\usepackage[a4paper,left=24mm,right=24mm,top=24mm,bottom=25mm,headheight=15pt]{geometry}
\usepackage{amsmath,amssymb,amsthm,mathtools,bm,mathrsfs}
\usepackage{fontspec}
\IfFontExistsTF{TeX Gyre Termes}{\setmainfont{TeX Gyre Termes}}{\setmainfont{Liberation Serif}}
\IfFontExistsTF{TeX Gyre Heros}{\setsansfont{TeX Gyre Heros}}{\setsansfont{Liberation Sans}}
\IfFontExistsTF{Latin Modern Mono}{\setmonofont[Scale=0.86]{Latin Modern Mono}}{\setmonofont[Scale=0.86]{Liberation Mono}}
\usepackage{graphicx,booktabs,array,tabularx,longtable}
\usepackage{microtype,caption,subcaption,fancyhdr,placeins,needspace,etoolbox}
\usepackage[numbers,sort&compress]{natbib}
\usepackage[hidelinks,bookmarksnumbered]{hyperref}
\usepackage{xurl}
\setlength{\emergencystretch}{2em}
\setlength{\parindent}{1.3em}
\setlength{\parskip}{0.16em}
\setlength{\textfloatsep}{14pt plus 3pt minus 3pt}
\setlength{\floatsep}{12pt plus 2pt minus 2pt}
\setlength{\intextsep}{12pt plus 2pt minus 2pt}
\setlength{\tabcolsep}{5pt}
\renewcommand{\arraystretch}{1.15}
\renewcommand{\topfraction}{0.9}
\renewcommand{\bottomfraction}{0.8}
\renewcommand{\textfraction}{0.08}
\renewcommand{\floatpagefraction}{0.7}
\setcounter{topnumber}{2}\setcounter{bottomnumber}{2}\setcounter{totalnumber}{4}
\captionsetup{font=small,labelfont=bf,labelsep=period,skip=6pt}
\captionsetup*[sub]{font=small,labelformat=parens}
\pagestyle{fancy}\fancyhf{}
\fancyhead[L]{\small Resistive thermoviscous evolution of sheared magnetic tori}
\fancyfoot[C]{\thepage}
\renewcommand{\headrulewidth}{0.3pt}
\newtheorem{theorem}{Theorem}
\newtheorem{lemma}{Lemma}
\newtheorem{proposition}{Proposition}
\theoremstyle{remark}\newtheorem{remark}{Remark}
\numberwithin{equation}{section}
\allowdisplaybreaks[2]
\pretocmd{\section}{\Needspace{6\baselineskip}}{}{}
\pretocmd{\subsection}{\Needspace{5\baselineskip}}{}{}
\newcommand{\dd}{\mathrm{d}}
\newcommand{\vv}{\boldsymbol v}
\newcommand{\BB}{\boldsymbol B}
\newcommand{\JJ}{\boldsymbol J}
\newcommand{\Avec}{\boldsymbol A}
\newcommand{\EE}{\boldsymbol E}
\newcommand{\MM}{\boldsymbol M}
\newcommand{\xx}{\boldsymbol x}
\newcommand{\XX}{\boldsymbol X}
\newcommand{\nn}{\boldsymbol n}
\newcommand{\xii}{\boldsymbol\xi}
\newcommand{\atest}{\boldsymbol a}
\newcommand{\ttau}{\boldsymbol\tau}
\newcommand{\etab}{\boldsymbol\eta}
\newcommand{\Id}{\boldsymbol I}
\newcommand{\qv}{\boldsymbol q}
\newcommand{\grad}{\nabla}
\newcommand{\diver}{\nabla\!\cdot}
\newcommand{\curl}{\nabla\!\times}
\newcommand{\Dt}{\mathrm D_t}
\newcommand{\Cv}{\mathcal C_{\vv}}
\newcommand{\st}{\mathfrak s}
\newcommand{\sym}{\operatorname{sym}}
\newcommand{\skw}{\operatorname{skw}}
\newcommand{\sech}{\operatorname{sech}}
\newcommand{\suppref}{Supplementary Material}

\usepackage{float}
\newcommand{\repositoryurl}{\url{https://github.com/hl-ren/pdel-lean-verification}}
\newcommand{\repositoryarchive}{\href{https://github.com/hl-ren/pdel-lean-verification/releases/tag/v1.0.1}{verification release v1.0.1}}

% END INCLUDED FILE

\hypersetup{pdftitle={A path-dependent entropic Lagrangian formulation for exact thermoviscous evolution of sheared magnetic tori with anisotropic resistivity},pdfauthor={Huilong Ren}}
\title{\bfseries A path-dependent entropic Lagrangian formulation for exact thermoviscous evolution of sheared magnetic tori with anisotropic resistivity}
\author{Huilong Ren\\[0.3em]
\small College of Civil Engineering, Tongji University, Shanghai 200092, China\\
\small \texttt{hlren@tongji.edu.cn}}
\date{}
\begin{document}
\maketitle
\begin{abstract}
Extending a three-dimensional magnetostatic equilibrium to resistive evolution requires the induction, momentum and thermal equations to remain compatible. This paper constructs exact thermoviscous evolutions of a non-axisymmetric magnetic torus with nonzero magnetic shear within a path-dependent entropic Lagrangian (PDEL) formulation. Starting from a known sheared toroidal equilibrium, a symmetric positive-definite resistivity tensor is constructed so that the resistive differential expression preserves the seed field. The resulting homogeneous expansion is coupled to an ideal gas and positive bulk viscosity without an imposed body force or volumetric heat supply. For a one-temperature gas with heat-capacity ratio $5/3$, the temperature-dependent resistivity has a $T^{-2/3}$ prefactor, and all evolving fields are elementary functions of physical time. A two-temperature extension uses a $T_e^{-3/2}$ prefactor and positive interspecies heat exchange on an explicitly bounded state interval. Material magnetic flux decays while the inherited rotational-transform profile is retained. Local total energy, moving-boundary power and transported dissipation histories close exactly. A seed-specific calculation excludes positive scalar resistivity from the same shape-preserving homologous trajectory, and the two-temperature exchange coefficient is shown to have a nonremovable endpoint pole. Symbolic identities, independent Cartesian differentiation and surface--volume quadrature verify the construction. The results provide passive inverse-constitutive realizations and exact multiphysics benchmarks; the prescribed anisotropic tensor and exchange law are not a classical collisional plasma closure.
\end{abstract}
\noindent\textbf{Keywords:} path-dependent entropic Lagrangian; resistive magnetohydrodynamics; exact solution; anisotropic resistivity; magnetic shear; energy balance.
% BEGIN INCLUDED FILE: sections/01_introduction.tex
\section{Introduction}\label{sec:intro}
Magnetohydrodynamic (MHD) models connect the geometry of a magnetic confinement region to force balance, magnetic-flux transport and the evolution of thermal energy. Classical toroidal equilibrium theory established the importance of magnetic surfaces, rotational transform and the constraints imposed by three-dimensional geometry \citep{Kruskal1958,Grad1967,Helander2014}. An equilibrium and a thermodynamically consistent evolution answer different questions: the former balances pressure and magnetic force, whereas the latter must also account for resistive flux loss, irreversible heating and power exchanged with the boundary. For a compressible conducting medium, a pressure field that satisfies instantaneous force balance need not satisfy an equation of state and a temperature equation simultaneously.

Analytic equilibria and exact dynamic solutions provide complementary ways to examine these requirements. Landreman \citep{Landreman} constructed explicit Cartesian non-axisymmetric equilibria with nested toroidal surfaces, including a family with nonconstant rotational transform. Issan et al.\ \citep{Issan} extended the corresponding families by parameters that remove stellarator symmetry. In a different mathematical construction, G\'omez-Serrano et al.\ \citep{Gomez2026} established smooth equilibria on embedded solid tori with prescribed discrete symmetry groups. These studies address existence and geometry, rather than the complete dissipative thermal evolution studied here. Exact shearing-box solutions with viscosity, resistivity and cooling, and nonlinear helical-wave constructions, illustrate how dynamic solutions can instead be selected by special compatibility conditions \citep{Lesaffre2007,Zhu2014}. Montani et al.\ \citep{Montani} make the magnetic--thermal compatibility issue explicit for resistive diffusion of toroidal equilibria. Together, these results motivate an exact benchmark that retains a genuinely three-dimensional toroidal geometry while satisfying all of the coupled balance equations.

The transport model is an essential part of such a benchmark. The classical Spitzer--H\"arm and Braginskii theories relate electrical and thermal transport to collisional plasma physics \citep{Spitzer,Braginskii1965}. Generalizations include axisymmetric equilibria with anisotropic resistivity and toroidal flow, and nonlinear resistive laws obtained from collisional momentum transfer \citep{Poulipoulis,Lingam}. These works also distinguish two different objectives: deriving transport coefficients from microscopic physics, and finding a passive constitutive realization of a specified continuum trajectory. The present paper pursues the second objective. Its resistivity tensor, bulk-viscosity texture and interspecies exchange law are constructed once from the reference solution and then transported with the material. Positivity and thermal compatibility are enforced, but the resulting coefficients are not identified with a classical collisional closure.

Variational formulations provide an organizing structure for this coupling. Newcomb's MHD action and Euler--Poincar\'e reduction establish the relation between material motion, advected quantities and Eulerian variations in the reversible setting \citep{Newcomb1962,Holm1998}. Variational nonequilibrium thermodynamics, port-Hamiltonian formulations and metriplectic descriptions extend energetic organization to irreversible processes and boundary power \citep{GayBalmaz2017a,GayBalmaz2017b,PortHamiltonian,Materassi2012,Coquinot2020}. For MHD specifically, Carlier and Campos-Pinto \citep{Carlier} develop a visco-resistive Hall formulation, while structure-preserving finite element methods enforce magnetic constraints and coupled energy or entropy balances at the discrete level \citep{Gawlik2021,Gawlik2025}. Exact fields with independently evaluable thermal and boundary powers are useful tests for these methods; a small magnetic residual alone does not check the complete multiphysics implementation.

The path-dependent entropic Lagrangian (PDEL) formulation introduces accumulated process powers into a state--history functional \citep{RenEntropy,Ren}. Its terminal differential distinguishes arbitrary weak field variations from the evaluation of the same constitutive action on realized rates; the endpoint and weak power-port structure is developed in \citet{RenSpending}. The submitted Eulerian formulation treats thermoviscous flow with species diffusion and reaction \citep{RenEulerian}, while the multi-channel formulation organizes simultaneous power conversion among interacting irreversible processes \citep{RenMultiChannel}. These studies supply the material-to-spatial history construction and the common donor--receiver bookkeeping used below. In the present retained-field MHD model, magnetic energy is included explicitly. One resistive response then supplies the vector Ohm equation, its mechanical Lorentz action and the Joule receiver. PDEL is used here as a particular common-power variational construction, not as the first variational treatment of dissipative MHD.

The constructive step is elementary but restrictive. If a reference vector potential has a strictly positive scalar product with the current, a symmetric positive-definite tensor can map that current to the potential. The curl of the resulting resistive field is proportional to the reference magnetic field. Homogeneous expansion then permits an explicit time dependence, provided density, viscosity and temperature are chosen together to satisfy momentum and thermal balance. We obtain a one-temperature branch with a $T^{-2/3}$ resistivity prefactor and a finite-interval two-temperature branch with a $T_e^{-3/2}$ prefactor. Both lose material magnetic flux while retaining the reference rotational-transform profile. A local scalar-resistivity obstruction and a nonremovable exchange endpoint quantify limitations of the construction. Parameter studies distinguish changes in geometry, the dissipative clock and the heat allocation; they do not imply stability or improved confinement.

The remainder of the paper is outlined as follows. Section~\ref{sec:review} reviews the PDEL operations through a thermo-viscoelastic example and their Eulerian and magnetic extensions. Section~\ref{sec:pdel} develops the present MHD functional and governing equations. Section~\ref{sec:realization} introduces the toroidal seed and its passive resistive realization. Sections~\ref{sec:oneT} and \ref{sec:twoT} establish the one- and two-temperature solutions, and Section~\ref{sec:energy} derives their energy and history identities. Section~\ref{sec:verification} presents parameter-dependent results and verification. Section~\ref{sec:discussion} summarizes the conclusions. Appendices~\ref{app:geometry}--\ref{app:review} collect the geometric calculation, constitutive restrictions and a detailed first-variation calculation for the introductory example.

% END INCLUDED FILE

% BEGIN INCLUDED FILE: sections/02_pdel_review.tex
\section{Review of the PDEL construction}\label{sec:review}
\subsection{A thermo-viscoelastic continuum and its explicit power histories}
The small-strain thermo-viscoelastic example fixes the state and history operations before their Eulerian extension \citep{RenEntropy,RenSpending,RenEulerian,RenMultiChannel}. Let $\Omega_0$ be a fixed reference body, $\boldsymbol u(\XX,t)$ its displacement, $\theta>0$ its temperature, and $s$ its entropy per reference volume. Set $\boldsymbol\varepsilon=\sym\grad_X\boldsymbol u$ and use the reference-volume Helmholtz density
\begin{equation}
\begin{aligned}
\psi(\boldsymbol\varepsilon,\theta)
={}&\tfrac12\boldsymbol\varepsilon:\mathbb C:\boldsymbol\varepsilon
-(\theta-\theta_0)\boldsymbol\beta_T:\boldsymbol\varepsilon\\
&+\rho_0c_v\left[(\theta-\theta_0)-\theta\log(\theta/\theta_0)\right].
\end{aligned}
\label{eq:reviewfree}
\end{equation}
Here $\rho_0c_v>0$ is the reference volumetric heat capacity, $\mathbb C$ is the elastic tensor, and $\boldsymbol\beta_T$ is the thermoelastic stress-coupling tensor. These coefficients are independent of temperature in this example. The reversible stress, Kelvin--Voigt stress and referential heat flux are
\begin{equation}
\begin{aligned}
\boldsymbol\sigma_e&=\psi_{,\boldsymbol\varepsilon}
=\mathbb C:\boldsymbol\varepsilon-(\theta-\theta_0)\boldsymbol\beta_T,\\
\boldsymbol\sigma_v&=\mathbb V:\dot{\boldsymbol\varepsilon},\qquad
D_v=\boldsymbol\sigma_v:\dot{\boldsymbol\varepsilon}
=\dot{\boldsymbol\varepsilon}:\mathbb V:\dot{\boldsymbol\varepsilon}\ge0,\\
\boldsymbol Q&=-\boldsymbol K_0\grad_X\theta,
\qquad \boldsymbol K_0=\boldsymbol K_0^T\succeq0.
\end{aligned}
\label{eq:reviewconstitutive}
\end{equation}
The viscosity tensor $\mathbb V$ is positive semidefinite on symmetric strain rates. With no volumetric heat supply, the two history densities have the explicit integrands
\begin{equation}
\begin{aligned}
H_Q(\XX,t)&=\int_{t_0}^{t}\grad_X\!\cdot\boldsymbol Q(\XX,\tau)\dd\tau
=-\int_{t_0}^{t}\grad_X\!\cdot(\boldsymbol K_0\grad_X\theta)(\XX,\tau)\dd\tau,\\
H_{D_v}(\XX,t)&=\int_{t_0}^{t}
\boldsymbol\sigma_v(\XX,\tau):\dot{\boldsymbol\varepsilon}(\XX,\tau)\dd\tau.
\end{aligned}
\label{eq:reviewhistories}
\end{equation}
Thus $\boldsymbol Q$ is a heat-flux vector, $\grad_X\!\cdot\boldsymbol Q$ is signed outward thermal power per reference volume, and $D_v$ is local irreversible conversion. For prescribed body-force density $\boldsymbol b_0$ and traction $\overline{\boldsymbol t}$, the state--history functional is
\begin{equation}
\begin{aligned}
\mathscr A_{\rm tv}={}&\int_{t_0}^{t_1}\!\int_{\Omega_0}
\left[\tfrac12\rho_0|\dot{\boldsymbol u}|^2-\psi-\theta s
-H_Q+H_{D_v}+\boldsymbol b_0\cdot\boldsymbol u\right]\dd V_0\dd t\\
&+\int_{t_0}^{t_1}\!\int_{\Gamma_t}
\overline{\boldsymbol t}\cdot\boldsymbol u\,\dd S_0\dd t.
\end{aligned}
\label{eq:reviewaction}
\end{equation}
Entropy and temperature remain independent mixed-state variables until the thermal state variation is taken. The histories organize realized power exchange; they are not additional physical stored energies.

\subsection{Independent variations and continuum field equations}
Keep an arbitrary compatible displacement variation $\delta\boldsymbol u$, with
$\delta\boldsymbol\varepsilon=\sym\grad_X\delta\boldsymbol u$, and an independent temperature variation $\delta\theta$. The displacement variation vanishes at the time endpoints and on the prescribed-displacement boundary. The independent entropy selector is $\delta t_s$, with $\delta^{\rm ch}s=\dot s\,\delta t_s$. The locked-history operations are
\begin{equation}
\begin{aligned}
\delta^{\rm ch}H_Q&=(\grad_X\!\cdot\boldsymbol Q)\,\delta t_s,\\
\delta^{\rm ch}H_{D_v}
&=D_v\,\delta t_s-\boldsymbol\sigma_v:\delta\boldsymbol\varepsilon.
\end{aligned}
\label{eq:reviewrules}
\end{equation}
The receiver term is the terminal derivative of the committed integral; the donor term is the weak action of the same constitutive stress \citep{RenSpending,RenMultiChannel}. On the realized mechanical strain direction,
\begin{equation}
\delta^{\rm ch}\boldsymbol\varepsilon
=\dot{\boldsymbol\varepsilon}\,\delta t_m,
\qquad
\boldsymbol\sigma_v:\delta^{\rm ch}\boldsymbol\varepsilon
=D_v\,\delta t_m,
\qquad
\delta^{\rm ch}H_{D_v}=D_v(\delta t_s-\delta t_m).
\label{eq:reviewcalibration}
\end{equation}
This is a calibration of the power pairing, not a restriction of the arbitrary mechanical test. In particular, for spatially varying $\delta t_m$, the ordinary strain variation is formed before this evaluation; it is not replaced by $\sym\grad_X(\dot{\boldsymbol u}\,\delta t_m)$.

Applying the state and channel operations to \eqref{eq:reviewaction} and integrating by parts gives
\begin{equation}
\begin{aligned}
\mathfrak D\mathscr A_{\rm tv}={}&-\int_{t_0}^{t_1}\!\int_{\Omega_0}
\Big\{\big[\rho_0\ddot{\boldsymbol u}
-\grad_X\!\cdot(\boldsymbol\sigma_e+\boldsymbol\sigma_v)-\boldsymbol b_0\big]
\cdot\delta\boldsymbol u\\
&\qquad +(\psi_{,\theta}+s)\delta\theta
+\big[\theta\dot s+\grad_X\!\cdot\boldsymbol Q-D_v\big]\delta t_s\Big\}\dd V_0\dd t\\
&+\int_{t_0}^{t_1}\!\int_{\Gamma_t}
\big[\overline{\boldsymbol t}-(\boldsymbol\sigma_e+\boldsymbol\sigma_v)\boldsymbol N\big]
\cdot\delta\boldsymbol u\,\dd S_0\dd t.
\end{aligned}
\label{eq:reviewfirstvariation}
\end{equation}
The independent variations recover
\begin{equation}
\begin{aligned}
\rho_0\ddot{\boldsymbol u}
&=\grad_X\!\cdot(\boldsymbol\sigma_e+\boldsymbol\sigma_v)+\boldsymbol b_0,\\
s=-\psi_{,\theta}
&=\boldsymbol\beta_T:\boldsymbol\varepsilon+\rho_0c_v\log(\theta/\theta_0),\\
\theta\dot s+\grad_X\!\cdot\boldsymbol Q&=D_v,
\qquad (\boldsymbol\sigma_e+\boldsymbol\sigma_v)\boldsymbol N
=\overline{\boldsymbol t}\quad\text{on }\Gamma_t.
\end{aligned}
\label{eq:reviewequations}
\end{equation}
Consequently,
\begin{equation}
\rho_0c_v\dot\theta
+\theta\boldsymbol\beta_T:\dot{\boldsymbol\varepsilon}
-\grad_X\!\cdot(\boldsymbol K_0\grad_X\theta)
=\dot{\boldsymbol\varepsilon}:\mathbb V:\dot{\boldsymbol\varepsilon}.
\label{eq:reviewtemperature}
\end{equation}
For $e_0=\psi+\theta s$, the same stress power cancels between mechanical and internal energy:
\begin{equation}
\partial_t\left(\tfrac12\rho_0|\dot{\boldsymbol u}|^2+e_0\right)
+\grad_X\!\cdot\left[\boldsymbol Q
-(\boldsymbol\sigma_e+\boldsymbol\sigma_v)^T\dot{\boldsymbol u}\right]
=\boldsymbol b_0\cdot\dot{\boldsymbol u}.
\label{eq:reviewenergy}
\end{equation}
Appendix~\ref{app:review} gives the pre-integration-by-parts expression, boundary budget and entropy-flux form. Related nonequilibrium variational theories also organize stress power and entropy \citep{GayBalmaz2017a,GayBalmaz2017b}; here their association is specified by the PDEL terminal-history and weak-power rules.

Layer~2 of the accompanying Lean development treats a homogeneous single-degree-of-freedom reduction of this continuum example. It does not certify the spatial integration by parts or boundary traction in \eqref{eq:reviewfirstvariation}. Its proof source, and the separate Eulerian MHD Layer~3 source, have the build status reported in Section~\ref{sec:formal}.

\subsection{Eulerian transport and multi-channel power conversion}
The submitted Eulerian formulation extends the preceding material construction to thermoviscous flow with species transport and reaction \citep{RenEulerian}. Two transport operators are needed from the outset. For a scalar field $f$, the material derivative follows its value along a moving parcel; for a density $w$ measured per current volume, the conservative rate also accounts for local volume change:
\begin{equation}
\Dt f:=\partial_t f+\vv\cdot\grad f,\qquad
\Cv w:=\partial_t w+\diver(w\vv)
=\Dt w+w\diver\vv.
\label{eq:reviewrates}
\end{equation}
The material derivative acts componentwise on vectors. The calligraphic $\Cv$ is a transport operator, distinct from the specific heat $c_v$. Thus $T$ is differentiated by $\Dt$, whereas $\rho$, $\st$ and the current-volume power histories use $\Cv$. For example, $\Cv\rho=0$ expresses mass conservation; it does not imply $\Dt\rho=0$ in a compressible flow.

The channel labels identify the participating physical processes. The mechanical channel $m$ carries stress work, the entropy channel $s$ receives the associated irreversible power, and the electromagnetic channel $e$ carries field work when the field is retained. Their realized-path selectors are $\delta t_m$, $\delta t_s$ and $\delta t_e$. These selectors calibrate the respective powers; they do not replace arbitrary displacement, temperature or vector-potential variations. In a two-temperature model, independent thermal selectors $\delta t_{s_e}$ and $\delta t_{s_i}$ distinguish the electron and ion receivers. Species channels are not needed in the nonreacting specialization considered here.

For the single-temperature fluid without a magnetic field, the state--history functional is
\begin{equation}
\mathscr A_E=\int\!\int_{\Omega(t)}
\left[\tfrac12\rho|\vv|^2-\Phi(\rho,T)-T\st
-H_{\mathrm{div}\,q}+H_\nu\right]\dd V\dd t.
\label{eq:revieweuleraction}
\end{equation}
The explicit histories follow the material parcel $\xx=\xx(\XX,t)$, with Jacobian $\mathcal J>0$:
\begin{equation}
\begin{aligned}
H_{\mathrm{div}\,q}(\xx(\XX,t),t)
&=\frac1{\mathcal J(\XX,t)}\int_{t_0}^{t}
\mathcal J(\XX,\tau)\,\diver\qv(\xx(\XX,\tau),\tau)\dd\tau,\\
H_\nu(\xx(\XX,t),t)
&=\frac1{\mathcal J(\XX,t)}\int_{t_0}^{t}
\mathcal J(\XX,\tau)\,[\ttau:\grad\vv](\xx(\XX,\tau),\tau)\dd\tau.
\end{aligned}
\label{eq:revieweulerhistories}
\end{equation}
More generally, the push-forward of one current-volume power density $g$ is
\begin{equation}
H_g(\xx(\XX,t),t)
=\frac1{\mathcal J(\XX,t)}\int_{t_0}^{t}
\mathcal J(\XX,\tau)g(\xx(\XX,\tau),\tau)\dd\tau,
\qquad \mathcal C_{\vv}H_g=g.
\label{eq:reviewtransport}
\end{equation}
The constitutive choices are $\qv=-\boldsymbol K\grad T$ and, in the present bulk-viscous model, $\ttau=\zeta(\diver\vv)\Id$; hence $D_\nu=\ttau:\grad\vv=\zeta(\diver\vv)^2$. To display how the equations follow from \eqref{eq:revieweuleraction}, retain an arbitrary Eulerian displacement $\xii$ and the standard relabeling directions \citep{Holm1998,RenEulerian}:
\begin{equation}
\begin{aligned}
\delta_\xi\vv&=\partial_t\xii+(\vv\cdot\grad)\xii-(\xii\cdot\grad)\vv,&
\delta_\xi\rho&=-\diver(\rho\xii),\\
\delta_\xi T&=-\xii\cdot\grad T,&
\delta_\xi\st&=-\diver(\st\xii).
\end{aligned}
\label{eq:revieweulervariations}
\end{equation}
The temperature test $\delta_T T=\delta T$ is independent of this relabeling. The entropy channel has $\delta^{\rm ch}\st=\Cv\st\,\delta t_s$, and the two history operations are
\begin{equation}
\delta^{\rm ch}H_{\mathrm{div}\,q}=(\diver\qv)\delta t_s,
\qquad
\delta^{\rm ch}H_\nu=D_\nu\delta t_s-\ttau:\grad\xii.
\label{eq:revieweulerspending}
\end{equation}
The last term is the mechanical donor and the first term in the same viscous rule is its entropy receiver. Evaluating the already formed stress pairing on the physical rate gives $D_\nu\delta t_m$; no replacement of the arbitrary field $\xii$ is made inside a spatial gradient.

For compactly supported spatial tests and vanishing temporal endpoint displacement, integration by parts uses $\Cv\rho=0$ and
\[
\delta_\xi\!\int\!\int\tfrac12\rho|\vv|^2\dd V\dd t
=-\int\!\int\rho\Dt\vv\cdot\xii\,\dd V\dd t,
\qquad
-\rho\grad\Phi_\rho+\Phi_T\grad T=-\grad p,
\quad p:=\rho\Phi_\rho-\Phi.
\]
The pressure identity follows by differentiating $p$; the mixed $T\st$ terms cancel in the mechanical coefficient without first imposing thermal conjugacy. Adding the history terms yields the bulk first differential
\begin{equation}
\begin{aligned}
\mathfrak D\mathscr A_E=\int\!\int_{\Omega(t)}\Big\{&
[-\rho\Dt\vv-\grad p+\diver\ttau]\cdot\xii
-(\Phi_T+\st)\delta T\\
&+[-T\Cv\st-\diver\qv+D_\nu]\delta t_s\Big\}\dd V\dd t.
\end{aligned}
\label{eq:revieweulerfirstvariation}
\end{equation}
Since $\xii$, $\delta T$ and $\delta t_s$ are independent, stationarity gives
\begin{equation}
\rho\Dt\vv=-\grad p+\diver\ttau,\qquad
\st=-\Phi_T,\qquad
T\Cv\st+\diver\qv=D_\nu.
\label{eq:revieweulerbalances}
\end{equation}
The mechanical, conjugacy and entropy equations therefore arise as three coefficients of the same differential. Section~\ref{sec:pdel} retains these operations and adds independent magnetic storage and its resistive power pairing.

The submitted multi-channel formulation \citep{RenMultiChannel} addresses the additional case in which one realized power is transferred between several processes. For an energy-transfer rate $Q_{i\to e}$, its history contributes $Q_{i\to e}(\delta t_{s_e}-\delta t_{s_i})$: the receiving and donating thermal channels use the same quantity with opposite signs. Its total-energy contribution cancels, while its entropy contribution is $Q_{i\to e}(T_e^{-1}-T_i^{-1})$. This is the multi-channel structure used in Section~\ref{sec:twoT}; it does not prescribe a microscopic exchange coefficient.

\subsection{Retained electromagnetic storage and explicit Joule history}
A fixed conducting medium provides a retained-electromagnetic-field specialization of the same common-power construction. To identify the storage signs and history integrands explicitly, choose a compatible temporal-gauge potential with $\EE=-\partial_t\Avec$, $\BB=\curl\Avec$, constant $\epsilon_0,\mu_0,\sigma_c>0$, and conduction current $\JJ_c=\sigma_c\EE$. With $\qv=-\boldsymbol K\grad T$, set
\begin{equation}
\begin{aligned}
H_q(\xx,t)&=\int_{t_0}^{t}\diver\qv(\xx,\tau)\dd\tau,\\
H_J(\xx,t)&=\int_{t_0}^{t}\JJ_c(\xx,\tau)\cdot\EE(\xx,\tau)\dd\tau
=\int_{t_0}^{t}\sigma_c|\EE(\xx,\tau)|^2\dd\tau.
\end{aligned}
\label{eq:reviewEMhistories}
\end{equation}
The histories are ordinary fixed-position integrals only because the carrier medium is stationary. The corresponding functional is
\begin{equation}
\mathscr A_{\rm EM}=\int\!\int_\Omega
\left[\tfrac12\epsilon_0|\EE|^2-\tfrac1{2\mu_0}|\BB|^2
-\psi_{\rm th}(T)-T\st-H_q+H_J\right]\dd V\dd t,
\label{eq:reviewmaxwellaction}
\end{equation}
with the explicitly calibrated operations
\begin{equation}
\begin{aligned}
\delta^{\rm ch}H_q&=(\diver\qv)\delta t_s,\\
\delta^{\rm ch}H_J&=\sigma_c|\EE|^2\delta t_s+\JJ_c\cdot\delta\Avec,
\qquad \delta^{\rm ch}\st=\partial_t\st\,\delta t_s.
\end{aligned}
\label{eq:reviewEMrules}
\end{equation}
On the realized electromagnetic direction $\delta^{\rm ch}\Avec=\partial_t\Avec\,\delta t_e=-\EE\delta t_e$, the last weak term is $-D_J\delta t_e$. In the variational derivation, however, $\delta\Avec$ remains an arbitrary independent vector field. Its coefficient gives $\curl\boldsymbol H=\epsilon_0\partial_t\EE+\JJ_c$, with $\boldsymbol H=\BB/\mu_0$, while $\delta T$ and $\delta t_s$ give $\st=-\partial_T\psi_{\rm th}$ and $T\partial_t\st+\diver\qv=\sigma_c|\EE|^2$. Faraday's law and $\diver\BB=0$ follow from the potential representation; compatible initially neutral Gauss data are preserved for these constant coefficients.

Retained field storage and the electromagnetic boundary flux satisfy
\begin{equation}
\partial_t\left(\frac{\epsilon_0|\EE|^2}{2}
+\frac{|\BB|^2}{2\mu_0}\right)
+\diver(\EE\times\boldsymbol H)=-\JJ_c\cdot\EE.
\label{eq:reviewpoynting}
\end{equation}
The distinction between field storage, boundary power and internal conversion is consistent with port-based and dissipative variational descriptions \citep{PortHamiltonian,Carlier,Gawlik2025}. The magnetic-quasistatic limit used below neglects electric field storage but retains magnetic storage and an independent vector test. For a moving medium,
\begin{equation}
\JJ\cdot\EE=\vv\cdot(\JJ\times\BB)
+\JJ\cdot(\EE+\vv\times\BB).
\label{eq:reviewEMsplit}
\end{equation}
The relative electric-field power is the resistive receiver after constitutive closure; the Lorentz term is field--motion exchange. This separation, rather than treating all electromagnetic work as heat, connects the continuum and multi-channel review to the full MHD formulation that follows.

% END INCLUDED FILE

% BEGIN INCLUDED FILE: sections/02_formulation.tex
\section{PDEL formulation of thermoviscous resistive MHD}\label{sec:pdel}
\subsection{Unified state--history functional}
Let $\Omega(t)$ be a material domain generated by an orientation-preserving motion $\xx=\xx(\XX,t)$. All quantities below are consistently nondimensionalized, with magnetic permeability $\mu_0=1$. The independent thermal state consists of temperature $T>0$ and entropy density $\st$ per current volume. A representative PDEL functional is
\begin{equation}
\mathscr A=\int_{t_0}^{t_1}\!\int_{\Omega(t)}
\left[\frac12\rho|\vv|^2-\Phi(\rho,T)-T\st
-\frac12|\BB|^2-H_q+H_{\rm J}+H_\nu\right]\dd V\dd t.
\label{eq:action}
\end{equation}
The instantaneous terms represent kinetic, thermal and magnetic storage. This functional extends the reviewed state--history construction \citep{RenEntropy,RenSpending,RenEulerian,RenMultiChannel} to a retained induction field. The histories $H_q$, $H_{\rm J}$ and $H_\nu$ represent heat transport, resistive conversion and viscous conversion, respectively. Their meaning is fixed by the transport and differential rules below; their committed interior histories are not subjected to an ordinary variation.

Mass conservation is kinematic admissibility in the material-to-Eulerian description \citep{Newcomb1962,Holm1998,RenEulerian}:
\begin{equation}
\rho(\xx(\XX,t),t)\mathcal J(\XX,t)=\rho_{\rm ref}(\XX),\qquad
\mathcal J=\det\frac{\partial\xx}{\partial\XX}>0,
\qquad \vv\circ\xx=\partial_t\xx.
\label{eq:kinematics}
\end{equation}
Consequently $\partial_t\rho+\diver(\rho\vv)=0$. The material and conservative rates introduced in \eqref{eq:reviewrates} are, respectively,
\begin{equation}
\Dt=\partial_t+\vv\cdot\grad,\qquad
\Cv f=\partial_tf+\diver(f\vv).
\label{eq:rates}
\end{equation}
A history $H_D$ of power density $D$ is the spatial representation of a material accumulation:
\begin{equation}
H_D=\rho\Gamma_D,\qquad
\Dt\Gamma_D=\frac D\rho,\qquad
\Cv H_D=D.
\label{eq:historytransport}
\end{equation}
Equivalently, $\mathcal JH_D=\int_{t_0}^{t}\mathcal J(\XX,\tau)D(\xx(\XX,\tau),\tau)\dd\tau$ for zero initial history. A fixed-position time integral does not generally satisfy this identity.

For the ideal gas, using the mixed thermal state of PDEL \citep{RenEntropy,RenEulerian}, choose
\begin{equation}
\Phi(\rho,T)=\rho c_vT[1-\log(T/T_*)]
+\rho RT\log(\rho/\rho_*),\qquad
c_v=\frac{R}{\gamma-1},\quad \gamma>1.
\label{eq:freeenergy}
\end{equation}
Independent temperature variation subsequently yields $\st=-\Phi_T$, and hence
\begin{equation}
u=\Phi+T\st=\rho c_vT,\qquad
p=\rho\Phi_\rho-\Phi=\rho RT.
\label{eq:eos}
\end{equation}
The conjugacy is not substituted into \eqref{eq:action} before taking its mixed thermal differential. The specific entropy is $s=\st/\rho=c_v\log(T/T_*)-R\log(\rho/\rho_*)$, up to an arbitrary reference constant.

\subsection{Independent weak directions and common dissipative powers}
In a temporal-gauge chart, use $\BB=\curl\Avec$ and $\EE=-\partial_t\Avec$. The arbitrary vector-potential test $\atest=\delta\Avec$ is independent of the arbitrary fluid displacement $\xii$. The mechanical directions are the constrained Eulerian relabeling variations \citep{Holm1998,RenEulerian}; the magnetic test is kept independent, as required for a retained dynamical field \citep{Carlier,Gawlik2025}:
\begin{equation}
\begin{aligned}
\delta_\xi\vv&=\partial_t\xii+[\vv,\xii],&
\delta_\xi\rho&=-\diver(\rho\xii),\\
\delta_\xi\st&=-\diver(\st\xii),&
\delta_\xi T&=-\xii\cdot\grad T,
\end{aligned}
\label{eq:variations}
\end{equation}
where $[\vv,\xii]=(\vv\cdot\grad)\xii-(\xii\cdot\grad)\vv$. The displacement vanishes at the time endpoints. Compactly supported spatial tests establish the bulk equations; boundary data are specified in Section~\ref{sec:boundary}.

For a symmetric positive-definite resistivity tensor $\etab$, define
\begin{equation}
\boldsymbol r_A=\partial_t\Avec-\vv\times\BB,\qquad
\boldsymbol R_A=\etab^{-1}\boldsymbol r_A,\qquad
D_{\rm J}=\boldsymbol r_A\cdot\etab^{-1}\boldsymbol r_A.
\label{eq:relative}
\end{equation}
The associated weak power is
\begin{equation}
\mathcal W_{\rm J}[\xii,\atest]
=\int_{\Omega(t)}\boldsymbol R_A\cdot(\atest-\xii\times\BB)\dd V.
\label{eq:weakohm}
\end{equation}
When evaluated on the physical directions $\xii=\vv\delta t_m$ and $\atest=(\partial_t\Avec)\delta t_m$, the integrand becomes $D_{\rm J}\delta t_m$. This calibration does not restrict either arbitrary weak test.

For bulk viscosity and Fourier conduction, let
\begin{equation}
\ttau=\zeta(\diver\vv)\Id,\qquad
D_\nu=\ttau:\grad\vv=\zeta(\diver\vv)^2,\qquad
\qv=-\boldsymbol K\grad T,
\label{eq:viscousheat}
\end{equation}
with $\zeta>0$ and $\boldsymbol K=\boldsymbol K^T\succeq0$. The same stress is used in the weak mechanical pairing $\ttau:\grad\xii$ and in $D_\nu$. A nonnegative shear viscosity may also be retained, since its deviatoric strain rate vanishes on the exact solutions constructed below.

For clarity, each history in \eqref{eq:action} is the push-forward of its specified integrand, not an independent unspecified term. With all integrands below evaluated at $(\xx(\XX,\tau),\tau)$,
\begin{equation}
\begin{aligned}
H_q(\xx(\XX,t),t)
&=\frac1{\mathcal J(\XX,t)}\int_{t_0}^{t}\mathcal J(\XX,\tau)\,
\diver(-\boldsymbol K\grad T)\dd\tau,\\
H_{\rm J}(\xx(\XX,t),t)
&=\frac1{\mathcal J(\XX,t)}\int_{t_0}^{t}\mathcal J(\XX,\tau)\,
\boldsymbol r_A\cdot\etab^{-1}\boldsymbol r_A\,\dd\tau,\\
H_\nu(\xx(\XX,t),t)
&=\frac1{\mathcal J(\XX,t)}\int_{t_0}^{t}\mathcal J(\XX,\tau)\,
\zeta(\diver\vv)^2\dd\tau.
\end{aligned}
\label{eq:explicitMHDhistories}
\end{equation}
After the vector-potential equation has established $\boldsymbol R_A=-\JJ$, the Joule-history integrand is equivalently $\JJ\cdot\etab\JJ$. This equivalence is a consequence of stationarity; the pre-stationarity definition in \eqref{eq:explicitMHDhistories} does not presuppose Ohm's law.


The PDEL history rules apply the common-power weak lifting reviewed in Section~\ref{sec:review} \citep{RenSpending,RenEulerian,RenMultiChannel}:
\begin{equation}
\begin{aligned}
\delta^{\rm ch}H_{\rm J}
&=D_{\rm J}\delta t_s-\boldsymbol R_A\cdot(\atest-\xii\times\BB),\\
\delta^{\rm ch}H_\nu&=D_\nu\delta t_s-\ttau:\grad\xii,\\
\delta^{\rm ch}H_q&=(\diver\qv)\delta t_s,
\end{aligned}
\label{eq:spending}
\end{equation}
with $\Cv H_q=\diver\qv$, $\Cv H_{\rm J}=D_{\rm J}$ and $\Cv H_\nu=D_\nu$. The entropy direction is $\delta^{\rm ch}\st=\Cv\st\,\delta t_s$, while the additional temperature variation $\delta T$ is arbitrary. Equations~\eqref{eq:variations} and \eqref{eq:spending} define distinct components of the mixed differential. In particular, the constitutive operator is evaluated at the current state; no additional variation of its past values is performed inside a locked history.

Ordinary variation of current storage, followed by the stated channel rules, gives the bulk expression
\begin{equation}
\begin{aligned}
\mathfrak D\mathscr A=\int\!\int_{\Omega(t)}\big\{&
[-\rho\Dt\vv-\grad p+\diver\ttau+\BB\times\boldsymbol R_A]\cdot\xii\\
&-(\JJ+\boldsymbol R_A)\cdot\atest-(\Phi_T+\st)\delta T\\
&+[-T\Cv\st-\diver\qv+D_{\rm J}+D_\nu]\delta t_s\big\}\dd V\dd t.
\end{aligned}
\label{eq:jointdifferential}
\end{equation}
Here $\JJ=\curl\BB$. The kinetic variation uses \eqref{eq:kinematics}; without mass admissibility its mechanical coefficient would additionally contain $-\vv[\partial_t\rho+\diver(\rho\vv)]$. The mixed thermal state contributes $-\grad p$ before elimination of entropy conjugacy. These observations keep the kinematic and thermal assumptions explicit.

Arbitrariness of the vector-potential test gives $\boldsymbol R_A=-\JJ$, so
\begin{equation}
\EE'=\EE+\vv\times\BB=\etab\JJ,
\qquad D_{\rm J}=\JJ\cdot\etab\JJ>0\quad(\JJ\ne0).
\label{eq:ohm}
\end{equation}
The mechanical and thermal tests give momentum balance and the entropy equation. Together with Faraday's identity and the equation of state, the resulting system is
\begin{subequations}\label{eq:system}
\begin{align}
\partial_t\rho+\diver(\rho\vv)&=0,\label{eq:mass}\\
\rho\Dt\vv&=-\grad p+\diver\ttau+\JJ\times\BB,\label{eq:mom}\\
\partial_t\BB&=\curl(\vv\times\BB)-\curl(\etab\JJ),
\quad\diver\BB=0,\label{eq:induction}\\
\rho c_v\Dt T+p\diver\vv&=D_{\rm J}+D_\nu-\diver\qv.
\label{eq:heat}
\end{align}
\end{subequations}
The complete vector Ohm equation is obtained from an independent vector test, not from a scalar contraction of an induction residual. Magnetic storage is retained in \eqref{eq:action}; neither Lorentz work nor Joule heat is added again as a separate prescribed source. The formulation applies to strictly positive resistivity and a compatible potential chart. The associated Eulerian/independent-potential derivation is represented by the Layer~3 Lean source package described in Section~\ref{sec:formal}. That source has not yet been accepted by a successful compilation in the present record; the analytical derivation here does not rely on a claim of completed formal certification. The ideal-resistivity limit requires a different constrained admissible class. Independent-potential dissipative formulations also occur in \citet{Carlier}; the specific structure here is the PDEL weak-power and history pairing.

\subsection{Material textures and constitutive specification}
The anisotropic tensor is represented by a symmetric material field $\MM$, and $\beta$ is a scalar viscosity parameter. One objective transport prescription is
\begin{equation}
\Dt\MM=\boldsymbol W\MM-\MM\boldsymbol W,\qquad
\Dt\beta=0,\qquad \boldsymbol W=\skw\grad\vv.
\label{eq:textures}
\end{equation}
These fields carry no stored energy in the model and are not independent sources of mechanical or thermal power. They specify a heterogeneous conducting continuum. On a homogeneous expansion, $\boldsymbol W=0$, and they remain constant along each material trajectory.

The one-temperature construction will use
\begin{equation}
\etab=\eta_*(T/T_*)^{-2/3}\MM,\qquad
\zeta=\beta\rho(T/T_*)^{1/3},\qquad \eta_*>0.
\label{eq:oneclosures}
\end{equation}
Their material textures are determined in Sections~\ref{sec:realization} and \ref{sec:oneT} and are subsequently fixed. The exponent $-2/3$ is selected by the coupled construction, whereas the two-temperature extension in Section~\ref{sec:twoT} uses a different thermal state to obtain a $-3/2$ electron-temperature prefactor.

% END INCLUDED FILE

% BEGIN INCLUDED FILE: sections/03_seed_realization.tex
\section{Toroidal seed and passive resistive realization}\label{sec:realization}
\subsection{Reference equilibrium and magnetic shear}
In reference Cartesian coordinates $\XX=(X,Y,Z)$, consider the sheared equilibrium of \citet{Landreman} with the constant complex-phase extension of \citet{Issan}. For $\epsilon,S,\lambda>0$ and $\tau_s\in\mathbb R$, write
\begin{equation}
\begin{aligned}
w&=X+iY,&K&=\bar w\sqrt{1+\epsilon/\bar w^2},\\
\Xi&=wK+\frac\pi2-S+i\tau_s,&W_0&=\frac{i\sin\Xi}{2K}.
\end{aligned}
\label{eq:complexseed}
\end{equation}
The square-root branch is fixed by $\Re(K/\bar w)>0$ on the retained domain. The magnetic field and seed pressure are
\begin{equation}
\begin{aligned}
B_{0x}+iB_{0y}&=e^{-i\lambda Z}W_0,\\
B_{0z}&=\lambda^{-1}\Re(e^{-i\lambda Z}\cos\Xi),\\
\psi&=\tfrac12\{\sin^2(\lambda Z)+\lambda^2B_{0z}^2\},
\qquad p_0=p_a-\psi/\lambda^2.
\end{aligned}
\label{eq:seedfield}
\end{equation}
They satisfy the magnetostatic identities
\begin{equation}
\grad_X\!\cdot\BB_0=0,\qquad
\JJ_0=\grad_X\!\times\BB_0,\qquad
\JJ_0\times\BB_0=\grad_Xp_0.
\label{eq:mhs}
\end{equation}
The corresponding current is
\begin{equation}
\begin{aligned}
J_{0x}&=\left[\lambda+\frac{4Y^2-2\epsilon}{\lambda}\right]B_{0x}
-\frac{4XY}{\lambda}B_{0y},\\
J_{0y}&=-\frac{4XY}{\lambda}B_{0x}
+\left[\lambda+\frac{4X^2+2\epsilon}{\lambda}\right]B_{0y},\\
J_{0z}&=\lambda B_{0z}.
\end{aligned}
\label{eq:current}
\end{equation}
In particular, a smooth global design potential on this domain is
\begin{equation}
\Avec_0=\lambda^{-1}(B_{0x},B_{0y},0),\qquad
\grad_X\!\times\Avec_0=\BB_0,
\qquad g=\Avec_0\cdot\JJ_0.
\label{eq:designpotential}
\end{equation}
The potential identity follows directly from the $Z$-derivatives of \eqref{eq:seedfield} and the last line of \eqref{eq:current}.

For $m=2\psi$, the conditions
\begin{equation}
0<m_b<1,\qquad S>\arcsin\sqrt{m_b}
\label{eq:regularity}
\end{equation}
define a smooth embedded solid torus $\Omega_0$ bounded by $m=m_b$. Appendix~\ref{app:geometry} gives its explicit surface map, inverse-coordinate argument and volume element. All derivatives in the governing equations are Cartesian Euclidean derivatives. The field is nonzero on the closed domain, including the magnetic axis, and $\BB_0\cdot\nn_0=0$ on its boundary.

The rotational transform is evaluated from the field direction rather than prescribed as an independent profile, following the field-line construction of \citet{Landreman,Issan} and the magnetic-surface convention discussed in \citet{Helander2014}. With the period integrals $\mathcal P(m)$ and $\mathcal R(m)$ defined in Appendix~\ref{app:geometry},
\begin{equation}
\iota(m)=\frac{\mathcal R(m)}{\mathcal P(m)},\qquad
\Phi_{t,0}(m)=\frac1{4\lambda}\int_0^m\mathcal P(s)\dd s.
\label{eq:transform}
\end{equation}
Writing $q_s=1/\iota$ and $\rho_\Phi=\sqrt{\Phi_{t,0}(m)/\Phi_{t,0}(m_b)}$, the flux-normalized shear is
\begin{equation}
\widehat s=\frac{\dd\log q_s}{\dd\log\rho_\Phi}
=-2\frac{\int_0^m\mathcal P(s)\dd s}{\mathcal P(m)}\frac{\iota_m}{\iota}.
\label{eq:shear}
\end{equation}
For $\epsilon=\tau_s=1/2$ and $S=3$, Appendix~\ref{app:geometry} establishes $\iota_\psi(0)/\iota(0)>31/288$. Nonzero shear near the axis therefore follows analytically, independently of numerical differentiation of the plotted profile.

Table~\ref{tab:parameters} specifies the reference geometry and the two thermal examples before their first graphical use. The parameter sweeps in Section~\ref{sec:verification} vary one declared parameter at a time; all quantities are dimensionless.

\begin{table}[H]
\centering
\caption{Parameters of the one- and two-temperature examples.}
\label{tab:parameters}
\begin{tabular}{@{}lll@{}}\toprule
Quantity & One temperature & Two temperatures\\\midrule
$(\epsilon,S,\tau_s,\lambda)$ & $(1/2,3,1/2,4)$ & Same\\
Outer surface $m_b$ & $0.9$ & $0.9$\\
Seed-pressure offset $p_a$ & $1$ & $2$\\
$(b_*,T_*,R,\eta_*)$ & $(1,1,1,0.005)$ & Same\\
$(\gamma,c_v,\kappa_0,U)$ & $(5/3,3/2,0.08,0.32)$ & Same\\
Electron fraction $r$ & Not applicable & $1/4$\\
Analytical interval & $a\ge1$ (finite time) & $1\le a<4^{6/5}$\\
Displayed thermal interval & $1\le a\le4$ & $1\le a\le4$\\\bottomrule
\end{tabular}
\end{table}

\begin{figure}[tbp]
\centering
\includegraphics[width=.67\linewidth]{figures/01_toroidal_geometry.pdf}
\caption{Reference non-axisymmetric magnetic torus and field lines for the parameters in Table~\ref{tab:parameters}. The geometry is the attributed sheared Landreman--Issan equilibrium. The displayed field lines are evaluated by monotone inverse quadrature on the actual Euclidean magnetic surfaces.}
\label{fig:geometry}
\end{figure}

\subsection{A positive-definite tensor realizing the resistive field}
The next observation supplies the constitutive link between the reference geometry and its resistive evolution.
\begin{lemma}[Passive tensor realization]\label{lem:tensor}
Let $\BB_0=\grad_X\!\times\Avec_0$, $\JJ_0=\grad_X\!\times\BB_0$ and $g=\Avec_0\cdot\JJ_0>0$. For $\ell>0$, define
\begin{equation}
\MM_0=\ell\frac{\Avec_0\otimes\Avec_0}{g}
+\Id-\frac{\JJ_0\otimes\JJ_0}{|\JJ_0|^2}.
\label{eq:tensor}
\end{equation}
Then $\MM_0$ is symmetric positive definite and satisfies
\begin{equation}
\MM_0\JJ_0=\ell\Avec_0,\qquad
\grad_X\!\times(\MM_0\JJ_0)=\ell\BB_0.
\label{eq:eigenidentity}
\end{equation}
If the fields are smooth on a compact domain and $g$ is bounded away from zero, $\MM_0$ is uniformly positive definite there.
\end{lemma}
\begin{proof}
For every vector $\boldsymbol z$,
\[
\boldsymbol z^T\MM_0\boldsymbol z
=\ell\frac{(\boldsymbol z\cdot\Avec_0)^2}{g}
+|\boldsymbol P_{\JJ_0^\perp}\boldsymbol z|^2.
\]
Both terms are nonnegative. Their simultaneous vanishing requires $\boldsymbol z$ to be parallel to $\JJ_0$ and perpendicular to $\Avec_0$, which implies $\boldsymbol z=0$ because $g>0$. Multiplication of \eqref{eq:tensor} by $\JJ_0$ proves the first relation in \eqref{eq:eigenidentity}; taking its curl proves the second. Uniform positivity follows by continuity and compactness.
\end{proof}

For the present seed, choose $\ell=\lambda^2$. Equation~\eqref{eq:current} implies
\begin{equation}
g\ge\left(1-\frac{2\epsilon}{\lambda^2}\right)|\BB_{0\perp}|^2>0
\qquad\text{when}\qquad\lambda^2>2\epsilon,
\label{eq:gpositive}
\end{equation}
where $\BB_{0\perp}=(B_{0x},B_{0y},0)$. The horizontal current matrix is the sum of $\lambda\Id_2$, a positive-semidefinite rank-one matrix and $\lambda^{-1}\operatorname{diag}(-2\epsilon,2\epsilon)$, which proves the bound. The explicit surface map gives $|\BB_{0\perp}|>0$ throughout the regular torus.

The tensor is designed once from the reference potential and then treated as fixed material data. Changing the design potential by a gauge gradient generally defines a different material tensor; it is not a transformation of the same constitutive model. Ordinary electromagnetic gauge freedom remains available for the evolving potential. Likewise, \eqref{eq:eigenidentity} is a differential identity with the compatible boundary traces used below, rather than a spectral assertion for a homogeneous perfectly conducting boundary problem.

Positive definiteness establishes passivity, not microscopic realizability. For example, conventional symmetric magnetic-field-aligned transport distinguishes parallel and perpendicular components \citep{Braginskii1965,Poulipoulis}, giving a tensor of the form
\begin{equation}
\etab=\eta_\parallel\widehat{\BB}\otimes\widehat{\BB}
+\eta_\perp(\Id-\widehat{\BB}\otimes\widehat{\BB}),
\label{eq:gyrotropic}
\end{equation}
and consequently maps $\JJ$ into $\operatorname{span}\{\JJ,\BB\}$. The present selected resistive field requires $\MM_0\JJ_0=\lambda^2\Avec_0$, but $\Avec_0\cdot(\JJ_0\times\BB_0)$ is not identically zero. At the example point with labels $(m,\theta,\zeta)=(0.4,0.43,0.71)$ it is approximately $-0.00565946$. Thus \eqref{eq:gyrotropic} cannot in general produce this particular relative electric field. An alternative electric-potential gradient or a different magnetic trajectory may alter that conclusion; it is not a nonexistence theorem for all field-aligned anisotropic models.

\subsection{Restriction on a scalar-resistivity realization}
The tensorial freedom in \eqref{eq:tensor} is essential to this particular shape-preserving construction.
\begin{proposition}[Seed-specific scalar obstruction]\label{prop:scalar}
For the seed \eqref{eq:seedfield} with \eqref{eq:regularity} and $\lambda^2>2\epsilon$, no positive scalar $r(\XX)$ satisfies
\begin{equation}
\grad_X\!\times(r\JJ_0)=\Lambda\BB_0
\label{eq:scalareigen}
\end{equation}
in a neighborhood of its outer magnetic-axis point for a spatial constant $\Lambda$.
\end{proposition}
\begin{proof}
Dotting \eqref{eq:scalareigen} with $\JJ_0$ fixes the only possible ratio $r/\Lambda$ wherever $\JJ_0\cdot(\grad_X\!\times\JJ_0)\ne0$. At the axis point $(x_a,0,0)$, both this denominator and $\JJ_0\cdot\BB_0$ are positive. Appendix~\ref{app:scalar} evaluates the full curl of the resulting ratio and obtains the strictly positive residual component
\begin{equation}
\left[\grad_X\!\times\left(\frac r\Lambda\JJ_0\right)-\BB_0\right]_z
=\frac{16x_a^3\cosh\tau_s}
{\lambda k_a(\lambda^2+4x_a^2)^2}>0,
\label{eq:scalarresidual}
\end{equation}
where $k_a=\sqrt{x_a^2+\epsilon}$ and $x_ak_a=S$. This contradicts \eqref{eq:scalareigen}.
\end{proof}
A homologous evolution that preserves this seed shape would require \eqref{eq:scalareigen} at each time. The proposition therefore excludes that scalar-resistive path even when the scalar coefficient is spatially nonuniform. It does not exclude other toroidal seeds, changing magnetic shapes or nonhomologous velocities.

% END INCLUDED FILE

% BEGIN INCLUDED FILE: sections/04_one_temperature.tex
\section{One-temperature exact evolution}\label{sec:oneT}
\subsection{Construction and complete field equations}
The supplementary package contains SymPy scripts for the exponent matching, stress and power identities, and material-history derivatives in this section, together with independent Cartesian residual checks. The separate Lean verification suite is available at \repositoryurl\ and archived as \repositoryarchive. Its actual build results, axiom dependencies and coverage limits are summarized in Section~\ref{sec:formal}.

\begin{theorem}[Resistive thermoviscous homologous evolution]\label{thm:oneT}
Let the reference seed satisfy \eqref{eq:mhs}, \eqref{eq:designpotential}, $g>0$ and $p_0>g/3$ on $\overline\Omega_0$. Choose $\gamma=5/3$, $b_*,T_*,\eta_*>0$ and set
\begin{equation}
\kappa_0=\lambda^2\eta_*,\qquad U=4\kappa_0,\qquad
a(t)=1+Ut,\qquad \XX=\xx/a(t).
\label{eq:scale}
\end{equation}
Define the reference density and viscosity texture by
\begin{equation}
\rho_0=\frac{b_*^2(12p_0-g)}{6c_vT_*},\qquad
\beta_0=\frac{b_*^2(3p_0-g)}{108\kappa_0\rho_0}.
\label{eq:materialdata}
\end{equation}
On $\Omega(t)=a(t)\Omega_0$, the fields
\begin{equation}
\begin{aligned}
\vv&=U\XX,&\rho&=a^{-3}\rho_0(\XX),\\
\BB&=b_*a^{-9/4}\BB_0(\XX),&\JJ&=b_*a^{-13/4}\JJ_0(\XX),\\
T&=T_*a^{-3/2},&p&=\rho RT,\\
\MM&=\MM_0(\XX),&\beta&=\beta_0(\XX)
\end{aligned}
\label{eq:onesolution}
\end{equation}
solve \eqref{eq:system} and \eqref{eq:textures} with \eqref{eq:oneclosures} for every finite $t\ge0$. Density, temperature, thermodynamic pressure, resistivity and bulk viscosity are positive, and both resistive and viscous powers are strictly positive.
\end{theorem}
\begin{proof}
The motion $\xx=a\XX$ has Jacobian $a^3$, so \eqref{eq:mass} follows directly. Its material acceleration is $\Dt\vv=\ddot a\XX=0$. The Eulerian time derivative and convective acceleration are individually nonzero and cancel. The velocity gradient is $(U/a)\Id$, with zero spin and zero deviatoric part; hence \eqref{eq:textures} is satisfied.

The viscosity and isotropic viscous stress $\ttau=\Pi\Id$ are
\begin{equation}
\zeta=\frac{b_*^2(3p_0-g)}{108\kappa_0}a^{-7/2},\qquad
\Pi=\frac{b_*^2(3p_0-g)}9a^{-9/2}.
\label{eq:viscoussolution}
\end{equation}
Since $R=2c_v/3$, the thermodynamic pressure obeys
\begin{equation}
p-\Pi=b_*^2a^{-9/2}p_0,\qquad
\JJ\times\BB=b_*^2a^{-11/2}\grad_Xp_0.
\label{eq:forcebalance}
\end{equation}
The three momentum components therefore vanish exactly.

From Lemma~\ref{lem:tensor}, the relative electric field is
\begin{equation}
\EE'=\etab\JJ=\kappa_0b_*a^{-9/4}\Avec_0(\XX),
\qquad
\curl\EE'=\frac{\kappa_0}{a}\BB.
\label{eq:relativefield}
\end{equation}
For homogeneous expansion the induction equation becomes
\[
\Dt\BB+\frac{2U}{a}\BB=-\curl\EE'.
\]
The left side evaluated from \eqref{eq:onesolution} equals $-U\BB/(4a)$, which is the right side because $U=4\kappa_0$. The solenoidal constraint follows from that of the seed.

The realized powers are
\begin{equation}
D_{\rm J}=\kappa_0b_*^2g\,a^{-11/2},\qquad
D_\nu=\frac43\kappa_0b_*^2(3p_0-g)a^{-11/2}.
\label{eq:powers}
\end{equation}
On the thermal side,
\begin{equation}
\rho c_v\Dt T+p\diver\vv
=\frac{\kappa_0b_*^2}{3}(12p_0-g)a^{-11/2}
=D_{\rm J}+D_\nu.
\label{eq:thermalclosure}
\end{equation}
Temperature is spatially uniform, so the Fourier flux vanishes for any retained positive-semidefinite conductivity tensor. The hypotheses give positivity of \eqref{eq:materialdata}, \eqref{eq:viscoussolution} and \eqref{eq:powers}.
\end{proof}

An explicit potential pair useful for independent verification is
\begin{equation}
\Avec=b_*a^{-5/4}\Avec_0(\XX),\qquad
\phi=\vv\cdot\Avec,\qquad
\EE=-\partial_t\Avec-\grad\phi.
\label{eq:potentialpair}
\end{equation}
It produces the same electric field as \eqref{eq:relativefield} and is gauge-equivalent to the temporal-gauge chart used in Section~\ref{sec:pdel}.

The evolving thermal pressure is
\begin{equation}
p=\frac{b_*^2}{9}(12p_0-g)a^{-9/2}.
\label{eq:thermalpressure}
\end{equation}
It is generally not constant on magnetic surfaces because $g$ is not a magnetic-surface function. The surface-aligned quantity is the effective pressure $p-\Pi$ in \eqref{eq:forcebalance}. Bulk viscosity is therefore required by the vector force balance as well as by the thermal equation; it cannot be removed while retaining the same fields.

\subsection{Constitutive exponent selection}
The exponents in Theorem~\ref{thm:oneT} follow from a restricted but explicit matching problem. Let $T=T_*a^{s_T}$, $\BB=b_*a^{\delta_B-2}\BB_0$, and $\etab=\eta_*(T/T_*)^\nu\MM_0$, with a linear scale factor and common pressure, viscous-stress and magnetic-stress exponents. Appendix~\ref{app:exponents} gives
\begin{equation}
s_T=\frac1\nu,\qquad
\delta_B=\frac{\nu+1}{2\nu},\qquad
U=-\frac{2\nu\kappa_0}{\nu+1},\qquad
q_\zeta=\nu+1,
\label{eq:exponents}
\end{equation}
where $\zeta=\beta_0\rho(T/T_*)^{q_\zeta}$ and $\nu\notin\{0,-1\}$. For $p_0,g>0$, a sufficiently large positive seed-pressure offset yields a passive expanding branch precisely in the large-offset window
\begin{equation}
\gamma>\frac43,\qquad -1<\nu<-\frac1{3(\gamma-1)}.
\label{eq:passivewindow}
\end{equation}
At $\gamma=5/3$, the choice $\nu=-2/3$ belongs to this interval and gives \eqref{eq:onesolution}. By contrast, substituting $\nu=-3/2$ into this same separation class produces contraction with a positive isotropic viscous stress, requiring a negative bulk-viscosity coefficient. The two-temperature construction below changes the thermal state space rather than changing the exponent in an otherwise unchanged one-temperature solution.

\subsection{Material flux loss and retained magnetic shear}
Under \eqref{eq:onesolution}, the direction of the field in material coordinates is unchanged. The transform and normalized shear on corresponding material surfaces are consequently retained, while magnetic flux through a material cut decreases:
\begin{equation}
\iota(m,t)=\iota(m,0),\qquad
\widehat s(m,t)=\widehat s(m,0),\qquad
\Phi_t(m,t)=b_*a^{-1/4}\Phi_{t,0}(m).
\label{eq:fluxshear}
\end{equation}
The area factor is $a^2$, whereas the field amplitude is $a^{-9/4}$. The additional $a^{-1/4}$ flux factor distinguishes the solution from ideal frozen-in expansion. Thus nonzero magnetic shear coexists with a nonzero resistive induction term, although the shear profile is not redistributed in time.

This gives a physical distinction that a static equilibrium cannot test: unchanged field-line winding does not imply conservation of magnetic flux through a material cut. Here the winding is preserved by the chosen shape-preserving resistive operator, while the field weakens faster than the ideal $a^{-2}$ dilution. The solution therefore separates resistive flux loss from rearrangement of the magnetic surfaces. It establishes neither magnetic reconnection nor improved confinement; it provides an exact nonideal evolution in which those effects are not conflated. Relative to the cited magnetostatic seeds \citep{Landreman,Issan}, the additional content is the simultaneous prediction of current, stress, temperature and irreversible power on the moving torus.

\begin{figure}[tbp]
\centering
\includegraphics[width=.67\linewidth]{figures/02_resistive_expansion.pdf}
\caption{Homologous expansion of the same material magnetic surface $m=0.7$. Blue solid, orange dashed and purple dash--dot meshes denote $a=1$, $1.5$ and $2$, respectively; the legend also gives their material-flux ratios $a^{-1/4}$. The mesh curves are surface-coordinate lines, not magnetic field-line trajectories. Equal Cartesian length scales are retained. The magnetic evolution is common to both thermal branches.}
\label{fig:expansion}
\end{figure}

% END INCLUDED FILE

% BEGIN INCLUDED FILE: sections/05_two_temperature.tex
\section{Two-temperature extension}\label{sec:twoT}
\subsection{Thermal channels and a positive exchange law}
A two-temperature state permits the same magnetic and mechanical evolution to use an electron-temperature resistivity prefactor $T_e^{-3/2}$. Let
\begin{equation}
c_e=c_i=\frac{c_v}{2},\qquad R_e=R_i=\frac R2,\qquad
\overline T=\frac{T_e+T_i}{2},\qquad
p=p_e+p_i=\rho R\overline T,
\label{eq:twostate}
\end{equation}
where $p_j=\rho R_jT_j$. The coefficients are per unit mixture mass. Both thermal components share the bulk velocity; separate electron and ion momenta or Hall drift are not part of this model.

Replace the thermal block in \eqref{eq:action} by
\[
\Phi_e(\rho,T_e)+T_e\st_e+\Phi_i(\rho,T_i)+T_i\st_i,
\]
with each free energy of the form \eqref{eq:freeenergy} using $c_j,R_j$. Independent thermal variations yield $\st_j=-\partial_{T_j}\Phi_j$. Assign Joule heat to electrons and viscous heat to ions, and use one exchange power with opposite signs in the two channels:
\begin{equation}
\begin{aligned}
\rho c_e\Dt T_e+p_e\diver\vv&=D_{\rm J}+Q_{i\to e},\\
\rho c_i\Dt T_i+p_i\diver\vv&=D_\nu-Q_{i\to e},\\
Q_{i\to e}&=G_{ie}(T_e,\XX)(T_i-T_e).
\end{aligned}
\label{eq:twoheat}
\end{equation}
Heat conduction can be retained for either species, but vanishes on the spatially uniform temperatures constructed here.

Choose $0<r<1$, put $T_{e0}=rT_*$ and define the material scalar
\begin{equation}
C_0(\XX)=\frac{4r}{9}(12p_0-g),\qquad
a_e=(T_e/T_{e0})^{-3/2}.
\label{eq:Czero}
\end{equation}
The constitutive laws are
\begin{equation}
\begin{aligned}
\etab&=\eta_*(T_e/T_{e0})^{-3/2}\MM,\\
\zeta&=\beta\rho\left(\frac{T_e+T_i}{2T_*}\right)^{1/3},\\
G_{ie}(T_e,\XX)&=\frac{\kappa_0b_*^2}{2T_*}a_e^{-4}
\frac{C_0(\XX)a_e^{5/6}-g(\XX)}{1-r a_e^{5/6}}.
\end{aligned}
\label{eq:twoclosures}
\end{equation}
Here $g$ and $C_0$ are fixed advected scalar textures. The exchange law depends on current electron temperature and material data, rather than laboratory time or a separately imposed heat source. Its specified passive state band is
\begin{equation}
\left(\frac{g}{C_0}\right)^{6/5}<a_e<r^{-6/5},
\label{eq:exchangeband}
\end{equation}
on which $G_{ie}>0$. The distinction between a temperature scaling and a collision-derived transport law is important \citep{Spitzer,Braginskii1965,Lingam}. Only the temperature prefactor in \eqref{eq:twoclosures} has the Spitzer exponent \citep{Spitzer}; neither $\MM$ nor $G_{ie}$ is derived from the classical collisional closure.

\subsection{Exact solution and entropy production}
\begin{theorem}[Finite-interval two-temperature realization]\label{thm:twoT}
Retain the seed, density, velocity and magnetic fields of Theorem~\ref{thm:oneT}, and suppose $p_0>g/3$ and $C_0>g>0$. Replace the temperature by
\begin{equation}
T_e=rT_*a^{-2/3},\qquad
T_i=2T_*a^{-3/2}-rT_*a^{-2/3}.
\label{eq:twotemperatures}
\end{equation}
For $1\le a<r^{-6/5}$, these fields satisfy mass, momentum, induction and both thermal equations \eqref{eq:twoheat} with \eqref{eq:twoclosures}. The temperatures, resistivity tensor, bulk viscosity and exchange coefficient are positive, with $T_i>T_e$ and $Q_{i\to e}>0$.
\end{theorem}
\begin{proof}
The mean temperature is $\overline T=T_*a^{-3/2}$. Total pressure, internal energy, viscosity and momentum therefore coincide with the one-temperature solution. In addition,
\[
(T_e/T_{e0})^{-3/2}=a,
\]
so the resistive field, induction equation and Joule power are unchanged. Substitution into the electron equation requires
\begin{equation}
Q_{i\to e}=\kappa_0b_*^2a^{-11/2}(C_0a^{5/6}-g).
\label{eq:exchangepower}
\end{equation}
Since $T_i-T_e=2T_*a^{-3/2}(1-ra^{5/6})$, \eqref{eq:twoclosures} gives exactly this power. The ion equation follows by subtracting the electron equation from the already established total thermal equation. The inequalities follow from $C_0>g$ and $a<r^{-6/5}$.
\end{proof}

With specific entropies $s_j=c_j\log T_j-R_j\log\rho+\text{constant}$, the total entropy balance is
\begin{equation}
\rho\Dt(s_e+s_i)
=\frac{D_{\rm J}}{T_e}+\frac{D_\nu}{T_i}
+G_{ie}\frac{(T_i-T_e)^2}{T_eT_i}>0.
\label{eq:twoentropy}
\end{equation}
The exchange conserves total thermal energy, but contributes positively to entropy production. The exact species entropy changes are
\begin{equation}
\begin{aligned}
s_e-s_{e0}&=\frac{2c_v}{3}\log a,\\
s_i-s_{i0}&=\frac{c_v}{2}\log\frac{T_i}{(2-r)T_*}+c_v\log a.
\end{aligned}
\label{eq:speciesentropy}
\end{equation}
The ion entropy attains its maximum at $a=(3/4r)^{6/5}$ when this point lies after the initial state. A subsequent decline in ion entropy is compatible with \eqref{eq:twoentropy}, because the ion balance contains outgoing exchange as well as viscous heat input. This branch distinguishes a decrease of one subsystem's entropy from negative total entropy production: energy transferred to the colder electrons produces entropy even while it lowers the ion contribution. It therefore supplies a coupled test of the exchange signs and temperature weights, rather than just two prescribed cooling curves. The exchange coefficient is inverse-designed from this trajectory, not a Coulomb-collision rate.

\subsection{Nonremovable endpoint of the exchange law}
The upper endpoint in Theorem~\ref{thm:twoT} is a constitutive restriction, not only a numerical stiffness threshold. Put $a_{\rm eq}=r^{-6/5}$. As $a\to a_{\rm eq}^{-}$, both temperatures approach the same positive value, but \eqref{eq:exchangepower} tends to
\begin{equation}
Q_{\rm eq}=\kappa_0b_*^2a_{\rm eq}^{-11/2}N_{\rm eq},\qquad
N_{\rm eq}=\frac{C_0}{r}-g=\frac{48p_0-13g}{9}>\frac g3>0.
\label{eq:Qendpoint}
\end{equation}
Furthermore,
\begin{equation}
T_i-T_e=\frac{5T_*}{3}a_{\rm eq}^{-5/2}(a_{\rm eq}-a)
+O((a_{\rm eq}-a)^2).
\label{eq:Tendpoint}
\end{equation}
Consequently,
\begin{equation}
G_{ie}\sim\frac{3\kappa_0b_*^2}{5T_*}
\frac{a_{\rm eq}^{-3}N_{\rm eq}}{a_{\rm eq}-a}.
\label{eq:Gpole}
\end{equation}
The pole is unavoidable under the stated assumptions. A bounded exchange coefficient cannot maintain this temperature path through equality while supplying the nonzero required exchange power. Regularizing the coefficient would change the exact thermal equations and would require a different temperature trajectory or power partition.

% END INCLUDED FILE

% BEGIN INCLUDED FILE: sections/06_energy_histories.tex
\section{Energy balance and transported power histories}\label{sec:energy}
The exact trajectories resolve three distinct questions: how magnetic flux is lost without changing the inherited winding, how a cooling material can generate entropy, and where its decreasing total energy goes. The first is answered by \eqref{eq:fluxshear}; the following balances identify the corresponding internal conversion and boundary transfer. Their value is that these quantities are fixed by the same fields and constitutive responses, rather than assigned independently after solving the magnetic problem.
\subsection{Local energy and the moving boundary}
For the one-temperature branch let $u=\rho c_vT$, and for the two-temperature branch let $u=\rho(c_eT_e+c_iT_i)$. Define
\begin{equation}
\mathcal E=\frac12\rho|\vv|^2+u+\frac12|\BB|^2.
\label{eq:energydef}
\end{equation}
Using the magnetic/thermal power split familiar from resistive MHD \citep{Carlier,Gawlik2025}, taking the scalar product of momentum with $\vv$, contracting induction with $\BB$, and adding the thermal equation gives, on both exact branches,
\begin{equation}
\partial_t\mathcal E+
\diver\left[(\tfrac12\rho|\vv|^2+u+p)\vv
-\ttau\vv+\EE\times\BB\right]=0.
\label{eq:localenergy}
\end{equation}
The heat flux is zero on these solutions. For general nonuniform temperatures, its contribution would be added to the energy flux. Resistive and viscous powers cancel between their donor equations and the thermal receiver. In the two-temperature case, exchange also cancels on adding the thermal channels, without cancelling its entropy production.

Let the total Cauchy stress be
\begin{equation}
\boldsymbol{\mathcal T}=-p\Id+\ttau+\BB\otimes\BB-\frac12|\BB|^2\Id.
\label{eq:totalstress}
\end{equation}
Combining Reynolds transport with the electromagnetic bulk--boundary power pairing \citep{PortHamiltonian,RenMultiChannel}, on $\Omega(t)$ one obtains
\begin{equation}
\frac{\dd}{\dd t}\int_{\Omega(t)}\mathcal E\dd V
=-\int_{\partial\Omega(t)}
[-\boldsymbol{\mathcal T}\vv+\EE'\times\BB]\cdot\nn\dd S.
\label{eq:movingenergy}
\end{equation}
The relative electric field $\EE'=\EE+\vv\times\BB$ is required alongside total traction work. This distinguishes energy closure from constancy of energy in a materially moving volume.

Define the reference integrals
\begin{equation}
I_p=\int_{\Omega_0}p_0\dd V,\qquad
I_g=\int_{\Omega_0}g\dd V,\qquad
I_B=\int_{\Omega_0}|\BB_0|^2\dd V.
\label{eq:integrals}
\end{equation}
Two reference identities make the boundary power explicit:
\begin{equation}
\begin{aligned}
\int_{\partial\Omega_0}(p_0+|\BB_0|^2/2)\XX\cdot\nn_0\dd S
&=3I_p+I_B/2,\\
\int_{\partial\Omega_0}(\Avec_0\times\BB_0)\cdot\nn_0\dd S
&=I_B-I_g.
\end{aligned}
\label{eq:referenceports}
\end{equation}
The first follows by testing the divergence-free seed Cauchy stress against $\XX$, using $\BB_0\cdot\nn_0=0$. The second follows from $\grad_X\cdot(\Avec_0\times\BB_0)=|\BB_0|^2-g$.

The outward mechanical and electromagnetic powers are therefore
\begin{equation}
\begin{aligned}
P_{\rm mech}^{\rm out}&=4\kappa_0b_*^2a^{-5/2}(3I_p+I_B/2),\\
P_{\rm em}^{\rm out}&=\kappa_0b_*^2a^{-5/2}(I_B-I_g).
\end{aligned}
\label{eq:outwardpowers}
\end{equation}
The total kinetic energy $K_0=\frac12\int_{\Omega_0}\rho_0U^2|\XX|^2\dd V$ is constant. Direct volume integration gives
\begin{equation}
E(t)=K_0+b_*^2a^{-3/2}
\left[\frac{12I_p-I_g}{6}+\frac{I_B}{2}\right],
\label{eq:integratedenergy}
\end{equation}
and hence
\begin{equation}
-\dot E=P_{\rm mech}^{\rm out}+P_{\rm em}^{\rm out}
=\kappa_0b_*^2a^{-5/2}(12I_p-I_g+3I_B).
\label{eq:globalenergy}
\end{equation}
The sign of the electromagnetic port is determined by $I_B-I_g$; it is not assumed to be outward. Consequently, decreasing magnetic energy cannot in general be identified with Joule heating alone: moving-boundary mechanical work and relative Poynting transfer also enter the budget. For the baseline geometry the electromagnetic contribution is inward, yet mechanical work is larger and the material-domain total energy decreases. This is a physically open energy balance, not a failure of conservation.

\subsection{Compatible initial and boundary data}\label{sec:boundary}
At $t=0$, prescribe the selected exact fields and the material textures. On the material wall, prescribe the velocity $\vv=U\XX$ and the tangential relative electric trace furnished by \eqref{eq:relativefield}. Its surface-curl compatibility maintains $\BB\cdot\nn=0$. The heat-flux data are homogeneous. The corresponding total traction is an output, not an additional independent condition imposed together with the velocity.

These boundary data support \eqref{eq:movingenergy}, but the wall is generally not a perfect conductor because its tangential relative electric field is nonzero. The reference surface is material and impermeable to mass, not energetically isolated. No exterior vacuum, coil circuit or material-wall response is included. The construction provides compatible exact initial and boundary data rather than a well-posedness theorem for arbitrary nearby data.

\subsection{Reference-volume histories and entropy}
Let a hat indicate power accumulated per reference volume. Direct integration of \eqref{eq:powers} with the material Jacobian $a^3$ gives
\begin{equation}
\widehat H_{\rm J}=\frac{b_*^2g}{6}(1-a^{-3/2}),\qquad
\widehat H_\nu=\frac{2b_*^2}{9}(3p_0-g)(1-a^{-3/2}).
\label{eq:histories}
\end{equation}
The current-volume densities are $H_j(\xx,t)=a^{-3}\widehat H_j(\XX,t)$ and satisfy $\Cv H_j=D_j$. Thus the accumulation and its dilatational dilution are distinguished as in the Eulerian PDEL history construction \citep{RenEulerian}, as required by \eqref{eq:historytransport}.

Writing $u_0=\rho_0c_vT_*=b_*^2(12p_0-g)/6$, the complete material internal-energy budget is
\begin{equation}
\begin{aligned}
\widehat H_{\rm J}+\widehat H_\nu&=\frac{u_0}{3}(1-a^{-3/2}),\\
\widehat W_p:=\int_0^t a^3p\diver\vv\dd\tau
&=\frac{4u_0}{3}(1-a^{-3/2}),\\
a^3u-u_0&=\widehat H_{\rm J}+\widehat H_\nu-\widehat W_p.
\end{aligned}
\label{eq:materialbudget}
\end{equation}
Equation~\eqref{eq:materialbudget} fixes a particularly transparent partition. With $\Delta=1-a^{-3/2}$, the material element receives irreversible heat $u_0\Delta/3$, performs pressure work $4u_0\Delta/3$, and therefore loses internal energy $u_0\Delta$. Thus pressure work is exactly four times the accumulated irreversible heat along this branch. At $a=2$, the internal energy per reference volume is $2^{-3/2}\simeq0.35355$ of its initial value even though both dissipative histories have increased. Cooling here measures net energy loss through expansion, not absence of irreversible conversion. The one-temperature specific entropy satisfies
\begin{equation}
s-s_0=\frac{c_v}{2}\log a,\qquad
\rho T\Dt s=D_{\rm J}+D_\nu>0.
\label{eq:oneentropy}
\end{equation}

For the two-temperature branch, the exchange history is
\begin{equation}
\widehat H_{ie}=b_*^2
\left[\frac{3C_0}{8}(1-a^{-2/3})-\frac g6(1-a^{-3/2})\right],
\qquad \partial_t\widehat H_{ie}=a^3Q_{i\to e}.
\label{eq:exchangehistory}
\end{equation}
It enters the electron receiver and ion donor with opposite signs. There is one physical exchange, not an additional total-energy source. Its entropy contribution is $Q_{i\to e}(T_e^{-1}-T_i^{-1})\ge0$.

Finally, the normalized local history $h=6\widehat H_{\rm J}/(b_*^2g)=1-a^{-3/2}$ reconstructs the one-temperature trajectory:
\begin{equation}
\begin{aligned}
a&=(1-h)^{-2/3},&T/T_*&=1-h,&\rho/\rho_0&=(1-h)^2,\\
\BB(a\XX,t)&=b_*(1-h)^{3/2}\BB_0(\XX),&&
\Phi_t/\Phi_{t,\mathrm{init}}&=(1-h)^{1/6}.
\end{aligned}
\label{eq:historycoordinate}
\end{equation}
This is an explicit state coordinate on the constructed solution manifold. Since it can be recovered from the instantaneous scale or temperature, it does not introduce an additional hereditary material law.

The distinction between reference and current histories also has an observable consequence. Along a material point, $\widehat H_{\rm J}$ increases monotonically, whereas its current-volume density is proportional to $a^{-3}(1-a^{-3/2})$ and reaches a maximum at $a=(3/2)^{2/3}$. Subsequent decrease of $H_{\rm J}$ reflects dilatational dilution of accumulated energy per current volume, not negative Joule power. This separates a transported PDEL history from a fixed-location temporal accumulation.

Taken together, these relations provide verification targets beyond a small induction residual: a coupled implementation must reproduce the effective pressure $p-\Pi$, the magnetic-flux decay, the cancellation of internal power transfers, the material-history transport and the boundary flux using the same constitutive data. They complement exact-equilibrium benchmarks and energy-consistent MHD discretizations \citep{Landreman,Issan,Gawlik2021,Gawlik2025}. Their role is to expose inconsistent coupling or missing transport terms; they do not themselves establish the accuracy or stability of a numerical solver.

% END INCLUDED FILE

% BEGIN INCLUDED FILE: sections/07_verification.tex
\section{Parameter-dependent results and verification}\label{sec:verification}
\subsection{Reference settings and admissible parameter families}
Table~\ref{tab:parameters}, introduced with the reference geometry, specifies the baseline examples. The two-temperature case changes the pressure offset and therefore the material density and viscosity, but retains the same magnetic surfaces. The parameters are nondimensional benchmark values, not a fit to a plasma experiment. Conservative bounds on the entire baseline reference domain are
\begin{equation}
1<G<5,\qquad X^2+Y^2<5,\qquad
|\BB_{0\perp}|^2<1,\qquad 0<g<\frac{37}{16}.
\label{eq:domainbounds}
\end{equation}
For $p_a=1$, $p_0>31/32$, hence
\begin{equation}
3p_0-g>\frac{19}{32},\qquad \rho_0>\frac{149}{144}>0.
\label{eq:passivitybounds}
\end{equation}
For $p_a=2$ and $r=1/4$, $p_0>63/32$ and
\begin{equation}
C_0-g=\frac{12p_0-10g}{9}>\frac1{18}>0.
\label{eq:twopassivity}
\end{equation}
These inequalities establish passivity independently of the sampled matrix eigenvalues. Table~\ref{tab:sweeps} lists the additional comparisons. Changing a parameter generally changes the inverse-designed material textures; the comparisons are between specified exact realizations, not perturbations of one fixed constitutive medium.
\begin{table}[H]
\centering\small
\caption{Parameter studies. Unlisted parameters retain Table~\ref{tab:parameters}; the phase study uses $p_a=2$ and the two-temperature studies use the two-temperature baseline.}
\label{tab:sweeps}
\begin{tabularx}{\linewidth}{@{}p{.17\linewidth}p{.32\linewidth}X@{}}\toprule
Parameter & Values or range & Isolated effect\\\midrule
Phase $\tau_s$ & $0,\ 0.25,\ 0.5,\ 0.75$ & Radial winding and shear; $\tau_s=0$ is the Landreman magnetic reference.\\
Resistive scale $\eta_*$ & $0.0025,\ 0.005,\ 0.01$ & Physical-time rate; $\beta_0$ and $U$ change consistently.\\
Pressure offset $p_a$ & $0.85\le p_a\le2.5$ & Density, bulk viscosity and heat partition; geometry fixed.\\
Electron fraction $r$ & $0.25,\ 0.375,\ 0.5$ & Temperature separation, ion entropy and interval endpoint.\\\bottomrule
\end{tabularx}
\end{table}
For each family, the code records the domain parameters and conservative positivity margins. In particular, the lower pressure offset still satisfies $3p_0-g>0$ using \eqref{eq:domainbounds}, and increasing $r$ above $1/4$ at $p_a=2$ preserves the initial exchange margin. Every two-temperature curve stops strictly below its own endpoint.

\subsection{Three-dimensional distributions of the exact fields}\label{sec:fields3d}
Figures~\ref{fig:spatialBJ}--\ref{fig:spatialpowers} evaluate the Cartesian fields directly on the solid torus, rather than interpreting a time-profile curve as a spatial distribution. A quarter-sector is omitted from the rendering to expose two radial cuts through the interior. The cuts are visualization surfaces only; they neither change $\Omega_0$ nor introduce new boundary conditions. All three coordinate axes use the same physical length scale, and the colorbars report the actual nondimensional field values. The displayed baseline is the one-temperature solution at $a=1$; all maps use the same geometry and viewpoint.

The magnetic and current magnitudes in Figure~\ref{fig:spatialBJ} exhibit substantial variation both along and across the magnetic surfaces. On the sampled full domain, $|\BB|$ ranges approximately from $0.1460$ to $0.4392$ and $|\JJ|$ from $0.9163$ to $2.6299$. These are sampled ranges, not certified extrema. The current is not a constant scalar multiple of the field: its horizontal components contain the spatially varying matrix in \eqref{eq:current}. Accordingly, magnetic intensity alone does not determine the Joule-power distribution. The latter also depends on the inverse-designed resistivity tensor through $D_{\rm J}=\JJ\cdot\etab\JJ$.
\begin{figure}[tbp]
\centering
\begin{subfigure}{.49\linewidth}
\includegraphics[width=\linewidth]{figures/26_3d_magnetic_magnitude.pdf}
\caption{Magnetic-field magnitude.}
\end{subfigure}\hfill
\begin{subfigure}{.49\linewidth}
\includegraphics[width=\linewidth]{figures/27_3d_current_magnitude.pdf}
\caption{Current-density magnitude.}
\end{subfigure}
\caption{Exact Cartesian magnetic and current fields on a cutaway view of the non-axisymmetric solid torus. A linear Turbo rainbow color scale encodes field values without illumination shading or height-based coloring. The exposed cuts reveal radial structure while the remaining exterior is the magnetic surface $m=m_b$.}
\label{fig:spatialBJ}
\end{figure}

Figure~\ref{fig:spatialrhop} resolves the weaker material-state heterogeneity. For this baseline, $RT=1$ at $a=1$, so the nondimensional pressure and density have the same numerical values, approximately $1.2738$--$1.3165$ over the sampled volume. Their similar appearance is an equation-of-state consequence, not a duplicated field. Neither field is generally constant on a magnetic surface, because its spatial factor contains $g$ as well as $p_0$. On a fixed magnetic surface, $p_0$ is constant while
\begin{equation}
p=\frac{b_*^2a^{-9/2}}9(12p_0-g),\qquad
\Pi=\frac{b_*^2a^{-9/2}}9(3p_0-g),\qquad
p-\Pi=b_*^2a^{-9/2}p_0.
\label{eq:fieldmap_pressure}
\end{equation}
Thus pressure and isotropic viscous stress have identical tangential variations, which cancel in the force-balancing effective pressure. The supplementary effective-pressure map is constant on the outer magnetic surface but varies across the exposed radial cuts. This distinguishes the present thermoviscous realization from interpreting $p_0$ itself as the evolving thermal pressure.
\begin{figure}[tbp]
\centering
\begin{subfigure}{.49\linewidth}
\includegraphics[width=\linewidth]{figures/28_3d_density.pdf}
\caption{Mass density.}
\end{subfigure}\hfill
\begin{subfigure}{.49\linewidth}
\includegraphics[width=\linewidth]{figures/29_3d_thermal_pressure.pdf}
\caption{Thermodynamic pressure.}
\end{subfigure}
\caption{Material-state distributions for the same exact field and time as Figure~\ref{fig:spatialBJ}. The colorbars intentionally resolve the small but nonzero variation. Equality of the numerical density and pressure values follows from $RT=1$ in this particular nondimensional state; it is not a general identity.}
\label{fig:spatialrhop}
\end{figure}

The two positive power densities in Figure~\ref{fig:spatialpowers} have different spatial patterns and markedly different absolute scales. At fixed $m$, $D_{\rm J}\propto g$ whereas $D_\nu\propto3p_0-g$; their tangential variations are consequently opposed. The sampled values are approximately $0.00159$--$0.02008$ for $D_{\rm J}$ and $0.28986$--$0.30888$ for $D_\nu$. The same Turbo palette is used with separate, explicitly labeled value ranges; similar colors in the two panels therefore do not imply similar heating magnitudes. Their integrated partition is quantified below by \eqref{eq:joulefraction}.
\begin{figure}[tbp]
\centering
\begin{subfigure}{.49\linewidth}
\includegraphics[width=\linewidth]{figures/30_3d_joule_power.pdf}
\caption{Joule conversion.}
\end{subfigure}\hfill
\begin{subfigure}{.49\linewidth}
\includegraphics[width=\linewidth]{figures/31_3d_viscous_power.pdf}
\caption{Bulk-viscous conversion.}
\end{subfigure}
\caption{Spatially heterogeneous but strictly positive dissipative powers. Their opposite variations along a fixed magnetic surface follow from their explicit $g$-dependence. The distinct colorbar scales preserve the fact that bulk-viscous dissipation dominates this baseline.}
\label{fig:spatialpowers}
\end{figure}

The geometrically corresponding maps at later times follow directly from the exact solution:
\begin{equation}
\begin{aligned}
|\BB|(a\XX,t)&=a^{-9/4}|\BB|(\XX,0),&
|\JJ|(a\XX,t)&=a^{-13/4}|\JJ|(\XX,0),\\
\rho(a\XX,t)&=a^{-3}\rho_0(\XX),&
p(a\XX,t)&=a^{-9/2}p(\XX,0),\\
D_{\rm J}(a\XX,t)&=a^{-11/2}D_{\rm J}(\XX,0),&
D_\nu(a\XX,t)&=a^{-11/2}D_\nu(\XX,0).
\end{aligned}
\label{eq:fieldmap_scalings}
\end{equation}
These relations preserve the material patterns but change their physical size and magnitude. In contrast, $T=T_*a^{-3/2}$ is spatially uniform at each time. The supplementary maps therefore use uniform temperature coloring and a common color and coordinate scale at $a=1$ and $a=2$. In particular, $\grad T=0$ and $\qv=0$ on this solution. Although the dissipative powers are nonuniform, their sum is proportional to the same material factor $12p_0-g$ as $\rho_0$, so the homogeneous temperature rate satisfies the local heat equation. Further three-dimensional maps show velocity, resistivity eigenvalue, bulk viscosity, reference-volume Joule history, entropy production and the two-temperature exchange fields.

\subsection{Magnetic geometry, phase dependence and the literature limit}
Figures~\ref{fig:geometry} and \ref{fig:expansion} show the actual reference torus and its homologous expansion. The field lines lie on embedded magnetic surfaces; their winding is not inferred from a visual projection. The phase study in Figure~\ref{fig:transform} evaluates the exact period formulas \eqref{eq:periods}. Only the one-dimensional integrals and the surface-to-flux coordinate are numerical. The derivative of the period ratio is evaluated by differentiating the integrands analytically, rather than by differentiating a fitted curve.

Increasing $\tau_s$ over the displayed range lowers the axis transform from approximately $5.92852$ at $\tau_s=0$ to $4.63571$ at $\tau_s=0.75$. It changes the radial profile more strongly than a uniform rescaling: the wall values decrease from $9.21831$ to $3.91271$, and an interior transform maximum appears in the displayed examples. For $\tau_s=0.5$, a zero of $\iota_m$ is located at $m\simeq0.823664$; for $\tau_s=0.75$, it moves to $m\simeq0.0384914$. The normalized shear changes sign at the same interior surfaces, with the opposite sign convention to $\iota_m$. These locations describe the evaluated parameter examples, not a proof of global uniqueness of extrema for all parameters.

The mechanism is visible in the separated periods: $\tau_s$ changes the hyperbolic factor in $\mathcal R(m)$, while $\mathcal P(m)$ is unchanged. Radial dependence prevents this change from reducing to a constant multiplier of $\iota$. By contrast, $\lambda$ changes the vertical length scale, magnetic flux magnitude and resistive rate but cancels from the transform ratio. This distinguishes a genuine change of winding from a change of geometry scale or traversal time, in the usual magnetic-surface interpretation \citep{Helander2014}.
\begin{figure}[htbp]
\centering
\begin{subfigure}{.48\linewidth}\includegraphics[width=\linewidth]{figures/14_phase_transform.pdf}
\caption{Rotational transform.}\end{subfigure}\hfill
\begin{subfigure}{.48\linewidth}\includegraphics[width=\linewidth]{figures/15_phase_shear.pdf}
\caption{Flux-normalized shear.}\end{subfigure}
\caption{Phase-dependent magnetic profiles from the exact period integrals, plotted against the flux radius of each configuration. The $\tau_s=0$ curve is the sheared Landreman reference field; nonzero phase gives the Issan extension. All configurations use $\epsilon=0.5$, $S=3$, $\lambda=4$ and $m_b=0.9$. The coordinate normalizes each configuration by its own boundary flux.}
\label{fig:transform}
\end{figure}

The comparison to existing analytic equilibria is exact at the magnetic level. Setting $\tau_s=0$ in \eqref{eq:complexseed}--\eqref{eq:seedfield} recovers the sheared family of \citet{Landreman}; nonzero $\tau_s$ gives the extension of \citet{Issan}. On the evolving domain,
\begin{equation}
\frac{\BB(a\XX,t)}{b_*a^{-9/4}}=\BB_0(\XX),\qquad
\frac{p-\Pi}{b_*^2a^{-9/2}}=p_0(\XX).
\label{eq:literaturerecovery}
\end{equation}
Thus the normalized field and effective pressure recover the cited equilibrium at every time. The physical thermal pressure is generally different, as required by \eqref{eq:thermalpressure}. No regular zero-resistivity limit of the \emph{complete} thermoviscous model is claimed: in this family $\beta_0\propto\eta_*^{-1}$, so setting $\eta_*=0$ would be a singular constitutive operation. The magnetic reference comparison must not be presented as convergence to a classical ideal-plasma model.

\subsection{Cooling, resistive flux loss and the expansion clock}
At a fixed material scale $a$, both thermal branches have the same magnetic evolution. Material flux is $2^{-1/4}\simeq0.84090$ of its initial value at $a=2$ and $4^{-1/4}\simeq0.70711$ at $a=4$, while the winding and normalized shear remain unchanged. Resistive flux loss therefore does not require a change in the winding of this specially constructed field. It is the amplitude relative to the material area, rather than the appearance of the magnetic surfaces, that detects nonideal induction.

Figure~\ref{fig:thermalenergy} separates cooling from entropy production. Temperature falls as $a^{-3/2}$, whereas the resistivity prefactor grows as $a$. Nevertheless, $D_{\rm J}$ and $D_\nu$ both decrease as $a^{-11/2}$: current dilution and declining strain rate dominate the increasing resistivity. The reference internal-energy budget \eqref{eq:materialbudget} shows that accumulated pressure work is four times the accumulated irreversible heating. The net cooling is therefore consistent with positive entropy production. Magnetic and thermal energies integrated over the expanding volume both scale as $a^{-3/2}$, whereas material kinetic energy remains constant. This coupled interpretation is essential when comparing with dissipative exact solutions or thermodynamically consistent schemes \citep{Lesaffre2007,Montani,Gawlik2025}; magnetic diffusion alone would not determine the cooling law.
\begin{figure}[htbp]
\centering
\begin{subfigure}{.48\linewidth}\includegraphics[width=\linewidth]{figures/21_thermal_mechanisms.pdf}
\caption{Thermal state and resistivity.}\end{subfigure}\hfill
\begin{subfigure}{.48\linewidth}\includegraphics[width=\linewidth]{figures/22_energy_scalings.pdf}
\caption{Energy and flux ratios.}\end{subfigure}
\caption{Mechanisms along the one-temperature trajectory. Increasing resistivity coexists with cooling, positive entropy production and decreasing integrated magnetic and internal energies. The magnetic and thermal energy curves coincide after normalization because they share the same time exponent, not because the two energies are equal.}
\label{fig:thermalenergy}
\end{figure}

Figure~\ref{fig:ratecomparison} varies $\eta_*$ at fixed physical time. Since $a=1+4\lambda^2\eta_*t$, larger $\eta_*$ accelerates expansion, cooling and material-flux loss. At $t=5$, the temperature ratios for $\eta_*=0.0025,0.005,0.01$ are approximately $0.41409,0.23853,0.11618$, and the flux ratios are $0.86334,0.78751,0.69853$. When plotted against $a$, these curves collapse onto the same exact state laws. This family varies the dissipative clock, not the material shear profile. Because $U$ and $\beta_0$ are also changed consistently, the comparison is not an experiment in which only electrical resistivity changes while all other material properties and initial velocities remain fixed.
\begin{figure}[htbp]
\centering
\begin{subfigure}{.48\linewidth}\includegraphics[width=\linewidth]{figures/16_resistivity_temperature.pdf}
\caption{Temperature.}\end{subfigure}\hfill
\begin{subfigure}{.48\linewidth}\includegraphics[width=\linewidth]{figures/17_resistivity_flux.pdf}
\caption{Material magnetic flux.}\end{subfigure}
\caption{Effect of the reference resistive scale on physical-time evolution. Each curve is an elementary exact solution with its corresponding passive material texture. The time coordinate is dimensionless; identical values of $a$, rather than identical times, give identical normalized states.}
\label{fig:ratecomparison}
\end{figure}

\subsection{Pressure offset and the partition of irreversible heat}
The volume-integrated Joule fraction has a closed expression because both dissipative powers share a time factor:
\begin{equation}
f_{\rm J}=\frac{\int D_{\rm J}\dd V}{\int(D_{\rm J}+D_\nu)\dd V}
=\frac{3I_g}{12I_p-I_g}.
\label{eq:joulefraction}
\end{equation}
For fixed geometry, let $V_0=|\Omega_0|$ and $I_\psi=\int_{\Omega_0}\psi\dd V$. Then $I_p=V_0p_a-I_\psi/\lambda^2$, while $I_g$ is independent of $p_a$. Therefore
\begin{equation}
\frac{\dd f_{\rm J}}{\dd p_a}
=-\frac{36I_gV_0}{(12I_p-I_g)^2}<0.
\label{eq:partitionderivative}
\end{equation}
Increasing the pressure offset increases the required density and bulk-viscosity texture but leaves the magnetic field, current and shear unchanged. Joule power is unchanged, while viscous heating increases. Figure~\ref{fig:partition} shows the resulting decrease of $f_{\rm J}$ from about $3.997\%$ at $p_a=0.85$ to $1.329\%$ at $p_a=2.5$; the baseline values are $3.380\%$ and $1.666\%$. These are fractions of irreversible conversion, not fractions of the net temperature change. Interspecies exchange is not a third total-energy source.
\begin{figure}[htbp]
\centering
\includegraphics[width=.49\linewidth]{figures/18_pressure_heat_partition.pdf}
\caption{Pressure-offset dependence of the heat partition at fixed magnetic geometry. The points mark the one- and two-temperature baseline offsets. The decline follows analytically from \eqref{eq:partitionderivative}; changing the offset changes the inverse-designed material, not its magnetic winding.}
\label{fig:partition}
\end{figure}
The dominance of bulk-viscous heating in these examples is a property of the selected realization. It prevents interpreting them as Joule-heating-dominated confinement scenarios. More conventional parallel/perpendicular resistive equilibria and collision-derived laws address different constitutive questions \citep{Poulipoulis,Lingam,Braginskii1965}; passivity of the present tensor alone does not identify it with those models.

\subsection{Two-temperature separation, entropy and finite validity}
The initial electron fraction changes the temperature separation without changing the mean-temperature or magnetic trajectory. In particular,
\begin{equation}
\frac{T_i}{T_e}=\frac{2}{r}a^{-5/6}-1,\qquad
a_{\rm eq}=r^{-6/5},\qquad
\frac{a_{s_i,\max}}{a_{\rm eq}}=(3/4)^{6/5}.
\label{eq:rcomparison}
\end{equation}
Increasing $r$ moves equality earlier. For $r=0.25,0.375,0.5$, the endpoint factors are respectively $5.27803,3.24461,2.29740$. The ion-entropy maxima occur at $3.73719,2.29740,1.62671$. Figure~\ref{fig:twothermal} shows how the temperature ratio approaches unity while ion entropy first rises and then decreases; Figure~\ref{fig:interval} displays the shrinking admissible interval. The corresponding total entropy remains increasing by \eqref{eq:twoentropy}. Cooling alone does not determine either species' entropy change: density dilution and the direction of interspecies transfer also enter.
\begin{figure}[htbp]
\centering
\begin{subfigure}{.48\linewidth}\includegraphics[width=\linewidth]{figures/19_two_temperature_ratio.pdf}
\caption{Ion/electron temperature ratio.}\end{subfigure}\hfill
\begin{subfigure}{.48\linewidth}\includegraphics[width=\linewidth]{figures/23_ion_entropy_sweep.pdf}
\caption{Ion specific entropy.}\end{subfigure}
\caption{Dependence on the initial electron fraction at $p_a=2$. Each curve ends at $0.98a_{\rm eq}$, strictly before its own constitutive singularity. The decline in ion entropy represents outgoing heat exchange; the total production remains positive.}
\label{fig:twothermal}
\end{figure}
\begin{figure}[htbp]
\centering
\includegraphics[width=.49\linewidth]{figures/20_two_temperature_interval.pdf}
\caption{Exact equality endpoint and ion-entropy maximum as functions of $r$. The upper curve is an excluded endpoint, not a stable final state or a numerically resolved continuation through the pole.}
\label{fig:interval}
\end{figure}
At equality the required exchange power is still strictly positive, as established by \eqref{eq:Qendpoint}--\eqref{eq:Gpole}. A bounded temperature-difference law could not maintain this path. The result therefore characterizes the finite scope of the inverse-designed two-temperature realization, rather than predicting a singular Coulomb equilibration process. The $T_e^{-3/2}$ scaling agrees with the named temperature prefactor, but the tensor and exchange coefficient remain different from classical collisional transport \citep{Spitzer,Braginskii1965}.

\subsection{Boundary energy transfer and physical interpretation}
For the one-temperature reference geometry, independent volume quadrature gives
\begin{equation}
\begin{aligned}
|\Omega_0|&\simeq3.1501412006,&I_p&\simeq3.0977071246,\\
I_g&\simeq0.4141467264,&I_B&\simeq0.3100467783.
\end{aligned}
\label{eq:quadraturenumbers}
\end{equation}
Since $I_B-I_g<0$, the relative Poynting power is inward. At $a=1$, the mechanical outward power is approximately $3.0234063$, the electromagnetic contribution is $-0.0083280$, and the net outward power is $3.0150783$, equal to $-\dot E$. The electromagnetic input is only $0.2762\%$ of the net outward power in magnitude; plotting both contributions on a single scale obscures its sign. Figure~\ref{fig:boundary} therefore uses separate panels with a common normalization.
\begin{figure}[htbp]
\centering
\begin{subfigure}{.48\linewidth}\includegraphics[width=\linewidth]{figures/24_boundary_mechanical.pdf}
\caption{Mechanical and net outward power.}\end{subfigure}\hfill
\begin{subfigure}{.48\linewidth}\includegraphics[width=\linewidth]{figures/25_boundary_electromagnetic.pdf}
\caption{Relative electromagnetic power.}\end{subfigure}
\caption{Moving-boundary energy balance, normalized by $P_{\rm net}(1)=-\dot E(0)$. Mechanical work removes energy, whereas the electromagnetic port supplies a smaller amount. The curves share an $a^{-5/2}$ factor. Energy closure is a complete port balance, not constant energy in an isolated torus.}
\label{fig:boundary}
\end{figure}
The pressure and magnetic stresses do work on the expanding boundary even though material acceleration is zero. A nonzero relative electric trace supplies electromagnetic power while material flux decreases. Neither feature contradicts positive bulk dissipation: boundary exchange and internal conversion are different terms. This distinction is the same one emphasized by bulk--boundary and inter-channel power formulations \citep{PortHamiltonian,RenMultiChannel}, and it is important for testing moving-mesh MHD implementations. The present exact boundary data do not replace a solved vacuum/coil/wall system.

\subsection{Independent verification and reproducibility}
Symbolic calculations check time exponents, powers, histories and the positive-minor expressions of the tensor. Cartesian automatic differentiation obtains current from the magnetic field independently of the explicit current formula. A separate high-precision implementation computes the electric field from \eqref{eq:potentialpair} before testing Ohm's law. It therefore does not make that law vanish by defining the field from the right-hand side. Table~\ref{tab:residuals} records the previously completed checks retained with this revision.
\begin{table}[htbp]
\centering\small
\caption{Recorded differential checks. Maxima are observed values at selected test points, not global error bounds or model-accuracy estimates.}
\label{tab:residuals}
\begin{tabularx}{\linewidth}{@{}p{.23\linewidth}Xr@{}}\toprule
Check & Method and coverage & Largest residual\\\midrule
Reference field & Cartesian derivatives at 48 points & $3.56\times10^{-15}$\\
One-temperature system & 192 space--time cases & $8.89\times10^{-16}$\\
Two-temperature system & 72 space--time cases & $2.67\times10^{-15}$\\
Independent implementation & 55-decimal arithmetic; 12 cases & $8.2\times10^{-56}$\\
Scalar obstruction & Symbolic third-order Cartesian expansion & Exact identity\\\bottomrule
\end{tabularx}
\end{table}
The physical role of these checks is specific: the momentum residual tests cancellation using $p-\Pi$, the thermal residual tests the balance between cooling, pressure work and heating, and the history residual tests transport of accumulated power rather than a fixed-position integral. The multiprecision values indicate agreement of independently implemented formulas, not physical accuracy to that number of digits. Symbolic and numerical agreement also does not establish stability.

For the additional parameter figures, the period integrals and their analytic derivatives were reevaluated at multiple angular resolutions and compared with adaptive quadrature. The largest refinement difference at the recorded control points was $1.8\times10^{-15}$; a separate finite-difference derivative check was below $1.7\times10^{-8}$. Independent volume quadratures for the pressure-offset analysis differed by less than $1.1\times10^{-11}$. Raw data, parameter values and the figure-generation script are included. These are quadrature checks of exact solution formulas, not PDE time-integration results. The original moving-boundary surface--volume errors decrease from $4.71\times10^{-10}$ to $9.2\times10^{-15}$ and $6.9\times10^{-16}$ with quadrature refinement. The checked evaluator additionally rejects nonfinite data, invalid toroidal branches, negative viscosity and states at or beyond the two-temperature endpoint.

\subsection{Lean proof layers and their current status}\label{sec:formal}
The formal-verification repository at \repositoryurl\ separates Lean proof checking from the symbolic and numerical residual evaluations supplied with the manuscript. Lean verifies precisely stated implications relative to their hypotheses and logical dependencies; a source file or a scan for placeholders alone is not a proof build \citep{LeanValidation}. The archived author-run Lean~4.19.0 evidence records fresh compilation and axiom auditing of 37 Layer~1, 66 Layer~2 and 217 Layer~3 theorem statements. The additional cross-layer bridge contributes 41 theorems, and the Revision3 modules contribute 46. Five validation-only interface checks are recorded separately. In total, 407 existing theorem statements and these five checks were compiled, and all 626 explicitly named declarations were audited with dependencies restricted to \texttt{Classical.choice}, \texttt{Quot.sound} and \texttt{propext}. The repository contains actual build logs, complete declaration types, source hashes and a mapping keyed by stable LaTeX labels; the versioned snapshot is \repositoryarchive.

Layer~1 covers finite-dimensional passivity and selected balance identities. Layer~2 is a single-degree-of-freedom thermo-viscoelastic specialization of the continuum example in Section~\ref{sec:review}. Layer~3 treats single-temperature, magnetic-quasistatic Eulerian PDEL--MHD with independent vector-potential variations and fixed-box integration. The bridge retains its stated normalization, reference-state and interface-port conditions. The Revision3 additions establish a nonmagnetic specialization and scalar results for the prescribed expansion branch: the cumulative pressure-work/irreversible-heat ratio and the current-volume Joule-history derivative, peak and monotonicity intervals.

These records do not certify the full Maxwell system, the two-temperature PDEs, the toroidal embedding, a complete three-dimensional realization of the homologous branch, or the moving-boundary energy theorem. The original theorem statements, assumptions and independent test spaces are preserved. The author-run evidence is distinct from subsequent GitHub-hosted CI results and from the symbolic and numerical checks reported above.

% END INCLUDED FILE

% BEGIN INCLUDED FILE: sections/08_discussion.tex
\section{Conclusions}\label{sec:discussion}
A PDEL formulation has been used to construct exact thermoviscous evolutions of an attributed non-axisymmetric magnetic torus with nonzero shear. A symmetric positive-definite, inverse-designed resistivity tensor preserves the seed's spatial structure while permitting material magnetic-flux loss. The coupled momentum and thermal equations determine the density, bulk viscosity and temperature scaling. The one-temperature solution is elementary for every finite forward time; the two-temperature extension realizes a $T_e^{-3/2}$ resistivity prefactor on a specified finite interval.

The solutions distinguish magnetic winding from flux conservation and cooling from negative entropy production. Their explicit power histories and moving-boundary balance show how irreversible conversion, expansion work and electromagnetic transfer combine without double counting. This makes them benchmarks for the simultaneous magnetic, mechanical and thermal equations, rather than magnetic-field profiles supplemented by independent heat sources. The PDEL contribution is the consistent use of one constitutive action in its weak field equation, realized power and thermal receiver.

The construction applies to the stated passive heterogeneous continuum, not to a classical collisional-plasma closure. The scalar-resistivity obstruction and the unavoidable two-temperature exchange pole delimit this realization. Magnetic-shear redistribution, perturbation stability and an exterior vacuum--coil--wall model remain outside the established results.

% END INCLUDED FILE

\section*{Data and code availability}
The analytical field evaluator, symbolic derivation scripts, parameter-sweep data and residual-verification programs are provided in the accompanying manuscript package. The separate Lean sources, pinned build instructions, axiom audits and source hashes are available at \repositoryurl, with the versioned verification archive identified by \repositoryarchive. Layer-specific coverage and actual validation status are specified in Section~\ref{sec:formal}. The analytical programs evaluate exact fields and residuals rather than time-step general MHD initial data.
% Funding, competing interests, CRediT roles
% and the final AI-use declaration must be confirmed by the author before submission.
\appendix
% BEGIN INCLUDED FILE: sections/appendices.tex
\section{Toroidal coordinates and the shear calculation}\label{app:geometry}
The Cartesian field is the sheared family of \citet{Landreman} with the complex-phase extension of \citet{Issan}. The parameterization below makes its regular magnetic surfaces and period integrals explicit for the present material-domain calculation; it is not a new static equilibrium family. For the seed \eqref{eq:seedfield}, set $m=2\psi$, $k=\sqrt m$ and define
\begin{equation}
\begin{aligned}
n(\zeta)&=\tau_s+\frac\epsilon2\sin2\zeta,&
\sigma(\theta)&=S+\arcsin(k\sin\theta),\\
h_g&=\sqrt{4\sigma^2+\epsilon^2},&
G&=\frac{h_g+\epsilon\cos2\zeta}{2},\\
f&=\sqrt{1-m\sin^2\theta},&
L&=\sqrt{\cosh^2n-m},\\
a_c&=\sqrt{(h_g-\epsilon)/2},&
b_c&=\sqrt{(h_g+\epsilon)/2}.
\end{aligned}
\label{eq:mapfunctions}
\end{equation}
The surface map is
\begin{equation}
\boldsymbol r(m,\theta,\zeta)=
\left(a_c\cos\zeta,\ b_c\sin\zeta,
\frac1\lambda\left[
\arctan\frac{k\cos\theta}{L}
-\arctan\frac{k\sin\theta\tanh n}{f}\right]\right).
\label{eq:surfacemap}
\end{equation}
Both angles are $2\pi$-periodic and both arctangents use the principal branch. Under \eqref{eq:regularity}, $\sigma\ge S-\arcsin\sqrt{m_b}>0$, so $a_c,b_c,G,f,L$ remain positive. The horizontal coordinates satisfy
\begin{equation}
a_cb_c=\sigma,\qquad b_c^2-a_c^2=\epsilon,\qquad
K=b_c\cos\zeta-i a_c\sin\zeta,\qquad
wK=\sigma+i\frac\epsilon2\sin2\zeta.
\label{eq:ellipticidentities}
\end{equation}
They are single-valued elliptic coordinates outside the excluded focal segment.

For a fixed horizontal point, introduce the disk variables
\[
u_d=\sin(\sigma-S)=k\sin\theta,
\qquad v_d=-\Im(e^{-i\lambda Z}\sin\Xi)=k\cos\theta.
\]
They satisfy $u_d^2+v_d^2=m$ and $e^{-i\lambda Z}\sin\Xi=L-ik\cos\theta$. On the selected branch, $\partial_Zv_d=\lambda L>0$, so the vertical inverse is unique; \eqref{eq:surfacemap} gives that inverse. It lies within $|\lambda Z|\le\arcsin\sqrt{m_b}<\pi/2$. All functions are smooth in $(u_d,v_d)$, including at $m=0$, where only the polar labels degenerate. The magnetic axis is the planar ellipse obtained by setting $m=0$.

Direct differentiation yields the volume element and contravariant field components
\begin{equation}
\dd V=\frac{G}{\lambda h_gfL}\dd\psi\dd\theta\dd\zeta,
\qquad B_0^\theta=\frac{h_gf}{2G}>0,
\qquad B_0^\zeta=\frac{L}{2G}>0.
\label{eq:jacobian}
\end{equation}
Together with the inverse argument, these relations establish the embedded solid torus and the nonvanishing field.

The separated field-line construction follows the analytic seed geometry \citep{Landreman,Issan}. The two-angle period ratio is consistent with the rotational-transform convention in toroidal confinement theory \citep{Kruskal1958,Helander2014}. The field-line equation separates:
\begin{equation}
\frac{\dd\theta}{h_g\sqrt{1-m\sin^2\theta}}
=\frac{\dd\zeta}{\sqrt{\cosh^2[\tau_s+\epsilon\sin2\zeta/2]-m}}.
\label{eq:fieldlineseparation}
\end{equation}
Its complete periods are
\begin{equation}
\begin{aligned}
\mathcal P(m)&=\int_0^{2\pi}
\frac{\dd\theta}
{\sqrt{4[S+\arcsin(\sqrt m\sin\theta)]^2+\epsilon^2}
\sqrt{1-m\sin^2\theta}},\\
\mathcal R(m)&=\int_0^{2\pi}
\frac{\dd\zeta}{\sqrt{\cosh^2[\tau_s+\epsilon\sin2\zeta/2]-m}}.
\end{aligned}
\label{eq:periods}
\end{equation}
The transform in \eqref{eq:transform} follows by counting the two angular periods. For a fixed $\zeta$ cut, multiplication of the Jacobian by $B_0^\zeta$ gives $\dd\Phi_{t,0}/\dd\psi=\mathcal P(2\psi)/(2\lambda)$, which proves the flux and shear formulas.

For the near-axis shear, write $N_j=\langle\sech^j n\rangle$ and $f_0(s)=(4s^2+\epsilon^2)^{-1/2}$. The period averages expand as
\[
\frac{\mathcal R(m)}{2\pi}=N_1+\frac m2N_3+O(m^2),\qquad
\frac{\mathcal P(m)}{2\pi}=f_0(S)+\frac m4[f_0(S)+f_0''(S)]+O(m^2).
\]
Hence
\begin{equation}
\frac{\iota_\psi(0)}{\iota(0)}
=\frac{N_3}{N_1}-\frac12-
\frac{2(8S^2-\epsilon^2)}{(4S^2+\epsilon^2)^2}.
\label{eq:axialshear}
\end{equation}
To bound the first term, the weight $\sech\sqrt z$ decreases with $z=n^2$. Its covariance with $z$ is therefore nonpositive, and $\tanh^2n\le n^2$ gives
\[
\frac{N_3}{N_1}=1-\mathbb E_{\sech n}(\tanh^2n)
\ge1-\tau_s^2-\frac{\epsilon^2}{8}.
\]
At $\epsilon=\tau_s=1/2$ and $S=3$, the first term is at least $23/32$ and the final geometric term is strictly less than $1/9$. Thus $\iota_\psi(0)/\iota(0)>23/32-1/2-1/9=31/288$, as used in the main text. The seed geometry and field-line structure are attributed to \citet{Landreman,Issan}; the present calculation verifies the parameter choice used for the dissipative realization.

\section{Axis calculation for the scalar-resistivity obstruction}\label{app:scalar}
Scalar and anisotropic resistive closures lead to different compatibility restrictions \citep{Poulipoulis,Montani}. The following calculation is specific to the attributed field \eqref{eq:seedfield} \citep{Landreman,Issan}; those references are the source of the magnetic field, not a prior proof of the obstruction derived here. Let $\boldsymbol K_0=\grad_X\!\times\JJ_0$. If \eqref{eq:scalareigen} held, its scalar product with $\JJ_0$ would imply
\begin{equation}
r\,\JJ_0\cdot\boldsymbol K_0=\Lambda\,\JJ_0\cdot\BB_0.
\label{eq:scalardot}
\end{equation}
At the outer magnetic-axis point $(x_a,0,0)$ define
\begin{equation}
x_a^2=\frac{\sqrt{4S^2+\epsilon^2}-\epsilon}{2},\qquad
k_a=\sqrt{x_a^2+\epsilon},\qquad
x_ak_a=S,\qquad D_a=\lambda^2+4x_a^2.
\label{eq:axisparameters}
\end{equation}
Both products on the right and left of \eqref{eq:scalardot}, excluding $r$ and $\Lambda$, are positive there. Thus $\Lambda>0$, and in a neighborhood the only candidate is
\begin{equation}
r_s:=\frac r\Lambda
=\frac{\JJ_0\cdot\BB_0}{\JJ_0\cdot\boldsymbol K_0}.
\label{eq:uniquescalar}
\end{equation}
Differentiating the full Cartesian expressions \eqref{eq:seedfield} and \eqref{eq:current}, and only then evaluating on the axis, gives
\begin{equation}
\begin{aligned}
r_s&=D_a^{-1},&
(r_s)_X&=-\frac{8x_a\lambda^2}{(D_a+2\epsilon)D_a^2},\\
J_{0y}&=\frac{D_a+2\epsilon}{2\lambda k_a}\cosh\tau_s,&
K_{0z}&=\frac{4x_a}{\lambda k_a}\cosh\tau_s,
\end{aligned}
\label{eq:axisderivatives}
\end{equation}
with $J_{0x}=B_{0z}=0$. Therefore
\begin{equation}
\left.\left[\grad_X\!\times(r_s\JJ_0)-\BB_0\right]\right|_{(x_a,0,0)}
=\left(-\frac{8x_a\sinh\tau_s}{D_a^2},\ 0,\
\frac{16x_a^3\cosh\tau_s}{\lambda k_aD_a^2}\right).
\label{eq:fullaxisresidual}
\end{equation}
The third component is strictly positive, proving Proposition~\ref{prop:scalar}.

The derivative in \eqref{eq:axisderivatives} cannot be obtained by differentiating the axis-only value $r_s=D_a^{-1}$ as if it were an off-axis formula. In particular, differentiation of $\JJ_0\cdot\boldsymbol K_0$ includes $J_{0z,X}K_{0z}$. This term is retained in the independent third-order Taylor calculation supplied with the implementation. For Table~\ref{tab:parameters}, the third residual component is approximately $0.01566689095$.

\section{Homologous exponent matching and passivity}\label{app:exponents}
Matching induction alone does not ensure thermal compatibility, as illustrated in other resistive exact-solution studies \citep{Lesaffre2007,Montani}. Here all exponents are determined from the particular PDEL-generated system in Section~\ref{sec:pdel}. Consider the restricted separation class
\begin{equation}
\begin{aligned}
a&=1+Ut,&T&=T_*a^{s_T},&
\rho&=a^{-3}\rho_0,\\
\BB&=b_*a^{\delta_B-2}\BB_0,&
\etab&=\eta_*(T/T_*)^\nu\MM_0,&
\zeta&=\beta_0\rho(T/T_*)^{q_\zeta},
\end{aligned}
\label{eq:generalansatz}
\end{equation}
with $U\ne0$ and nonzero resistive action. The symbol $\delta_B$ is a magnetic-amplitude exponent, distinct from the material viscosity texture $\beta_0$.

Using \eqref{eq:eigenidentity}, induction requires
\begin{equation}
\nu s_T=1,\qquad U\delta_B=-\kappa_0.
\label{eq:inductionmatching}
\end{equation}
Pressure and magnetic stress share an exponent if $s_T-3=2\delta_B-4$, whereas viscous stress shares it if $-4+q_\zeta s_T=2\delta_B-4$. Consequently, for $\nu\notin\{0,-1\}$,
\begin{equation}
s_T=\frac1\nu,\qquad
\delta_B=\frac{\nu+1}{2\nu},\qquad
U=-\frac{2\nu\kappa_0}{\nu+1},\qquad
q_\zeta=\nu+1.
\label{eq:matchingresult}
\end{equation}
Writing $\Pi=\Pi_0a^{s_T-3}$, momentum fixes
$\Pi_0=\rho_0RT_*-b_*^2p_0$. The thermal equation then gives
\begin{equation}
\begin{aligned}
\rho_0&=-\frac{b_*^2}{c_vT_*}
\left[\frac{\nu+1}{2}g+3\nu p_0\right],\\
\Pi_0&=-b_*^2\left[
\frac{(\gamma-1)(\nu+1)}2g+
\{1+3\nu(\gamma-1)\}p_0\right],\\
\beta_0&=\frac{\Pi_0}{3U\rho_0}.
\end{aligned}
\label{eq:generalmaterial}
\end{equation}
These formulas reproduce \eqref{eq:materialdata} at $\nu=-2/3$, $\gamma=5/3$.

For expansion, $U>0$ forces $-1<\nu<0$. With $p_0,g>0$, a sufficiently large positive pressure offset gives $\rho_0>0$, while $\Pi_0>0$ additionally requires $1+3\nu(\gamma-1)<0$. This yields \eqref{eq:passivewindow}; at its upper boundary the remaining $g$-term makes $\Pi_0<0$. The interval is nonempty only for $\gamma>4/3$.

At $\nu=-3/2$ and $\gamma=5/3$, \eqref{eq:matchingresult} instead gives $U=-6\kappa_0$, while \eqref{eq:generalmaterial} gives $\Pi_0=b_*^2(2p_0+g/6)>0$. Positive density would therefore require $\beta_0<0$. This calculation excludes that exponent from the present passive one-temperature homologous class, not from all possible thermoviscous evolutions. The two-temperature branch changes the pressure--temperature matching and so lies outside this one-temperature obstruction.


\section{Continuum thermo-viscoelastic first variation and energy balance}\label{app:review}
The reference-volume model of Section~\ref{sec:review} follows the thermo-mechanical PDEL state and power conventions \citep{RenEntropy,RenSpending,RenEulerian,RenMultiChannel}. This appendix displays the continuum calculation without replacing compatible displacement variations by scalar trajectory tests. Put $\boldsymbol\sigma=\boldsymbol\sigma_e+\boldsymbol\sigma_v$ and $\delta\boldsymbol\varepsilon=\sym\grad_X\delta\boldsymbol u$. The mixed state differential is
\begin{equation}
\widehat\delta(\psi+\theta s)
=\boldsymbol\sigma_e:\delta\boldsymbol\varepsilon
+(\psi_{,\theta}+s)\delta\theta
+\theta\dot s\,\delta t_s.
\label{eq:reviewmixedstate}
\end{equation}
Both histories in \eqref{eq:reviewhistories} have their past integrands locked. Combining \eqref{eq:reviewmixedstate} with \eqref{eq:reviewrules}, before integration by parts, gives
\begin{equation}
\begin{aligned}
\mathfrak D\mathscr A_{\rm tv}
={}&\int_{t_0}^{t_1}\!\int_{\Omega_0}
\Big[\rho_0\dot{\boldsymbol u}\cdot\delta\dot{\boldsymbol u}
-\boldsymbol\sigma:\sym\grad_X\delta\boldsymbol u
-(\psi_{,\theta}+s)\delta\theta\\
&\qquad-\big(\theta\dot s+\grad_X\!\cdot\boldsymbol Q-D_v\big)\delta t_s
+\boldsymbol b_0\cdot\delta\boldsymbol u\Big]\dd V_0\dd t\\
&+\int_{t_0}^{t_1}\!\int_{\Gamma_t}
\overline{\boldsymbol t}\cdot\delta\boldsymbol u\,\dd S_0\dd t.
\end{aligned}
\label{eq:reviewbeforeIBP}
\end{equation}
Time-endpoint vanishing of $\delta\boldsymbol u$ removes the kinetic endpoint term. Stress symmetry gives $\boldsymbol\sigma:\sym\grad_X\delta\boldsymbol u=\boldsymbol\sigma:\grad_X\delta\boldsymbol u$. The spatial divergence theorem yields
\begin{equation}
-\int_{\Omega_0}\boldsymbol\sigma:\grad_X\delta\boldsymbol u\,\dd V_0
=\int_{\Omega_0}(\grad_X\!\cdot\boldsymbol\sigma)\cdot\delta\boldsymbol u\,\dd V_0
-\int_{\partial\Omega_0}(\boldsymbol\sigma\boldsymbol N)\cdot\delta\boldsymbol u\,\dd S_0.
\label{eq:reviewstressIBP}
\end{equation}
These steps give \eqref{eq:reviewfirstvariation}, its bulk equations and the natural traction. No entropy conjugacy was imposed in evaluating the first variation.

After obtaining $s=-\psi_{,\theta}$, differentiation of $e_0=\psi+\theta s$ and use of the entropy balance give
\begin{equation}
\dot e_0
=\boldsymbol\sigma_e:\dot{\boldsymbol\varepsilon}+\theta\dot s
=\boldsymbol\sigma_e:\dot{\boldsymbol\varepsilon}
+\boldsymbol\sigma_v:\dot{\boldsymbol\varepsilon}
-\grad_X\!\cdot\boldsymbol Q.
\label{eq:reviewinternal}
\end{equation}
The kinetic energy rate follows independently from momentum:
\begin{equation}
\partial_t\big(\tfrac12\rho_0|\dot{\boldsymbol u}|^2\big)
=\grad_X\!\cdot(\boldsymbol\sigma^T\dot{\boldsymbol u})
-\boldsymbol\sigma:\dot{\boldsymbol\varepsilon}
+\boldsymbol b_0\cdot\dot{\boldsymbol u}.
\label{eq:reviewkinetic}
\end{equation}
Adding \eqref{eq:reviewinternal} and \eqref{eq:reviewkinetic} proves \eqref{eq:reviewenergy}, and integration over the fixed reference body gives
\begin{equation}
\begin{aligned}
\frac{\dd}{\dd t}\int_{\Omega_0}
\big(\tfrac12\rho_0|\dot{\boldsymbol u}|^2+e_0\big)\dd V_0
={}&\int_{\Omega_0}\boldsymbol b_0\cdot\dot{\boldsymbol u}\,\dd V_0\\
&+\int_{\partial\Omega_0}
(\boldsymbol\sigma\boldsymbol N)\cdot\dot{\boldsymbol u}\,\dd S_0
-\int_{\partial\Omega_0}\boldsymbol Q\cdot\boldsymbol N\,\dd S_0.
\end{aligned}
\label{eq:reviewboundarybudget}
\end{equation}
On $\Gamma_t$ the traction is prescribed; on a moving prescribed-displacement boundary the corresponding reaction work is retained. The histories do not add another storage term to this physical energy.

Finally, the entropy-flux form distinguishes thermal transport from production:
\begin{equation}
\dot s+\grad_X\!\cdot\left(\frac{\boldsymbol Q}{\theta}\right)
=\frac{D_v}{\theta}-\frac{\boldsymbol Q\cdot\grad_X\theta}{\theta^2}
=\frac{\dot{\boldsymbol\varepsilon}:\mathbb V:\dot{\boldsymbol\varepsilon}}{\theta}
+\frac{\grad_X\theta\cdot\boldsymbol K_0\grad_X\theta}{\theta^2}\ge0.
\label{eq:reviewentropyflux}
\end{equation}
The divergence of $\boldsymbol Q$ itself need not be nonnegative. This stress-power and heat-flux decomposition is consistent with other continuum nonequilibrium variational formulations \citep{GayBalmaz2017b}; its PDEL representation retains the same realized history in the weak donor and entropy receiver.

The accompanying Layer~2 Lean source treats the homogeneous scalar reduction, not the spatial boundary terms established here. Layer~3 addresses the Eulerian MHD variational structure in a distinct source package. Their current build status remains that reported in Section~\ref{sec:formal}; neither a continuum traction theorem nor the moving-torus energy theorem is inferred from the compiled finite-dimensional auxiliary identities.

% END INCLUDED FILE

\FloatBarrier
\begingroup\small
\setlength{\bibsep}{2.5pt}
\interlinepenalty=10000
\bibliographystyle{unsrtnat}
\bibliography{references}
\endgroup
\end{document}
